\newcommand{\HEADTITLE}{质量度量下的拉格朗日 Hodge–Dirac：完整报告}
\documentclass[10.5pt,a4paper]{article}
\usepackage[a4paper,margin=2.2cm,top=2.4cm,bottom=2.4cm]{geometry}
\usepackage{fontspec}
\usepackage{amsmath,amssymb,mathtools,bm,amsthm}
\setmainfont{Noto Serif CJK SC}[BoldFont={Noto Serif CJK SC Bold}]
\setsansfont{Noto Sans CJK SC}[BoldFont={Noto Sans CJK SC Bold}]
\XeTeXlinebreaklocale "zh"
\XeTeXlinebreakskip = 0pt plus 1pt minus 0.1pt
\linespread{1.2}
\usepackage{xcolor}\definecolor{accent}{HTML}{1F5FA8}\definecolor{quiet}{HTML}{5A5A5A}\definecolor{tint}{HTML}{EEF3FA}
\usepackage{titlesec}
\titleformat{\section}{\sffamily\bfseries\Large\color{accent}}{\thesection}{0.8em}{}
\titleformat{\subsection}{\sffamily\bfseries\large}{\thesubsection}{0.6em}{}
\usepackage{enumitem}\setlist{itemsep=0.2em,topsep=0.3em,leftmargin=1.6em}
\usepackage{array,booktabs,tabularx,makecell}\renewcommand{\arraystretch}{1.22}
\newcolumntype{L}{>{\raggedright\arraybackslash}X}
\usepackage{graphicx,caption}\captionsetup{font=small,labelfont=bf,labelsep=quad}
\renewcommand{\figurename}{图}\renewcommand{\tablename}{表}\renewcommand{\contentsname}{目录}
\usepackage{fancyhdr}\pagestyle{fancy}\fancyhf{}
\fancyhead[L]{\small\color{quiet}\HEADTITLE}
\fancyhead[R]{\small\color{quiet}2026 年 10 月}\fancyfoot[C]{\small\thepage}
\usepackage[colorlinks=true,linkcolor=accent,urlcolor=accent]{hyperref}
\theoremstyle{definition}\newtheorem{definition}{定义}[section]
\theoremstyle{plain}\newtheorem{theorem}[definition]{定理}\newtheorem{proposition}[definition]{命题}\newtheorem{corollary}[definition]{推论}
\theoremstyle{remark}\newtheorem{remark}[definition]{注}
\renewcommand{\proofname}{证明}
\newcommand{\re}{\operatorname{Re}}\newcommand{\im}{\operatorname{Im}}
\newcommand{\AFH}{AFRC$_{\mathrm H}$}
\newcommand{\status}[1]{\hfill{\small\color{accent}\textsf{[#1]}}}
\setlength{\parskip}{0.35em}\setlength{\parindent}{0pt}\tolerance=2000 \emergencystretch=3em


\setcounter{tocdepth}{2}
\begin{document}
\begin{center}
{\sffamily\bfseries\Huge 质量度量下的拉格朗日 Hodge–Dirac 计算}\\[0.6em]
{\sffamily\Large 只凭示踪点位置的可压缩物质输运与掺混：理论、验证与应用（完整报告）}\\[0.8em]
{\color{quiet}2026 年 10 月（v2：修正与验证版）}
\end{center}
\vspace{0.5em}
\noindent\fcolorbox{accent!40}{tint}{\begin{minipage}{0.965\linewidth}\vspace{0.4em}
\textbf{\sffamily 总摘要}\quad
本报告把拉格朗日轨迹网络上的工作整合为一套只用示踪点位置（加一个参考时刻的密度）的理论与算法，并给出全部检验。
\textbf{理论}：三点映射 $(p,q)$ 的有限步恒等式——复上链 $(d_1c)_t=2A\,\mathrm{Im}\,p+i\,\Delta A$ 同时给出环量（T4）与面积变化（T2 账本），位置形式的离散环量恒等式严格成立，由它可推出离散 Kelvin 定理；这些恒等式与度量无关，给出物质账本；以账本质量为度量的 Dirac 算子的 Laplacian 块，就是变密度扩散（$L_0=\delta_1d_0$，T8）与加权动态 Laplacian（T7）；Nanson 恒等式把界面掺混精确分解为切向拉伸与压缩（T9）。
\textbf{验证}：七个算例（双环流、Lamb–Oseen、Sod、柱面爆炸波、激波–重气泡、激波–等熵涡、单模 Richtmyer–Meshkov），参照为解析解、精确解与 DNS（含三档网格加密与粘性 DNS）。密度中位误差至多 1.3\%（v3 混合估计为 0.11\%–0.47\%），比 DTFE、SPH 准 1–2 个数量级；扩散的质量形式比体积形式准 3.8--10.6 倍，差距来自方程，在 $\rho D\propto T^{0.7}$ 下依然存在；在同一物质边界上，位置形式环量与 DNS 速度线积分的中位相差 0.03\%–0.25\%，与 Samtaney–Zabusky 模型相差 8\%；全部代数恒等式通过精确符号验证，账本在制造流映射上得到验证；RM 振幅与 DNS 相差 $\le1.2\%$，增长率与 Richtmyer 模型相差 3\%。
\textbf{稳健性}：噪声传播律（中位误差 $\approx0.74\,\sigma/h$）、自适应锚定、稀疏与断轨（逐轨账本，覆盖率 100\%）。
\textbf{边界}：多边形账本在光滑流中为一阶（张量账本为二阶）；无粘 DNS 卷起后的界面带不收敛（粘性 DNS 中收敛）；爆炸波薄壳中断轨后新轨迹的标定；三维与真实 PTV 数据尚未做。
\vspace{0.3em}\end{minipage}}

\vspace{0.6em}
\noindent\textbf{\sffamily 阅读指南}\quad 第一部分给出 T1–T9 的定稿理论与逐项验证；第二部分把成熟部分重组为“质量度量下的 Dirac 算子”，并补做 DNS 加密、基线、参数敏感性、实际数据检验与 RM 应用，含统一数值登记表；第三部分补齐粘性 DNS、自适应锚定、噪声传播律、动态 Laplacian、环量沉积模型比较与附件文献定位；第四部分是断轨下的逐轨账本；第五部分是场图；第六部分（v2）是修正、补证与验证；附录部分是后续工作（脊分类）。总结部分给出最终状态总表、论文计划与代码说明。各部分保留了当时的结论；若后文对前文有修订，以后文与总结部分的登记表为准。

\begin{table}[!htbp]\centering\small\caption{最终状态总表（以本报告全部结果为准）。}
\begin{tabularx}{\linewidth}{@{}l >{\raggedright\arraybackslash}p{0.27\linewidth} L l@{}}\toprule
组成 & 理论 & 关键结果 & 去向\\\midrule
T1 激波检测 & 单步面积台阶＋多窗口检测（签名确认）；合并块签名（充要） & 检出 0.96–1.00；弱激波和慢激波制造算例由失效变为精确 & 论文 A\\
T2 物质账本 & 逐点锚定、跨激波推迟重锚；复上链虚部 & 密度中位 $\le1.3\%$（混合估计 0.11–0.47\%）；平面激波处精确；弯曲激波误差随 $h/R_s$；比 DTFE/SPH 准 1–2 量级；噪声律；逐轨断轨账本 & 论文 A（核心）\\
T3 局部读数 & 复上链实部；Kabsch 对齐 & $\hat\omega$ 与参照相关 0.97–0.9999（去掉跨激波三角形） & 论文 A（局部量）\\
T3 全局分解、T5 & 谱分割、声学 Dirac、DESP & 能量占比非局部、不作物理读数 & 拓扑信号另文\\
T4 Kelvin & 位置形式严格成立；区域收支 & 同边界线积分中位相差 0.03–0.25\%；与 SZ 模型 8\% & 论文 A\\
T6 质量输运 & 质量 0‑链推前 & 通量 0.16\%–1.4\%；RM 重气体质量 $\pm4\%$ & 论文 A\\
T7 相干集 & 加权动态 Laplacian（账本质量） & Lamb–Oseen 涡核 Jaccard 0.11→0.82 & 论文 A\\
T8 扩散 & $L_0=\delta_1d_0$，质量度量、$\rho\kappa$ 电导；方差衰减无条件 & 质量形式 0.2\%–6\%；体积形式大 3.8--10.6 倍（来自方程） & 论文 A（核心）\\
T9 掺混 & Nanson 精确分解；链式张量；信息极限 & 精确算例 $\le0.007$；RM 压缩占 94\%；SB $t\le0.6$ 先切向缩短 & 论文 A（核心）\\
脊分类 & 增长指数 $\beta$、跳跃占比 $K$、面积占比 $W$ & 剪切层 100\%；FTLE 激波候选判对 99.8\%–100\% & 第三项工作\\\bottomrule
\end{tabularx}\end{table}

\tableofcontents

\clearpage\part{第一部分　T1–T9：定稿理论与逐项验证}


\noindent\fcolorbox{accent!40}{tint}{\begin{minipage}{0.965\linewidth}\vspace{0.3em}
\textbf{\sffamily 范围}\quad
本稿只收已经成型的 T1–T9（以及 T10 噪声检验），把散在四份报告里的理论与数值统一到一处，并逐项判断理论是否完备、是否得到验证。三项内容不在本稿：
\begin{enumerate}[leftmargin=1.4em,itemsep=0.05em]
\item 基于 \AFH{}/FTLE 的脊分类（双曲 LCS、剪切型脊、激波）是后续工作，另成一稿；
\item Dirac 算子的声学解释与 DESP 联合滤波没有定量检验，只在 T5 中作为备注；
\item 三维推广与激波—边界层相互作用（SBLI）的应用尚未做。
\end{enumerate}

\textbf{\sffamily 相对旧版的变化}\quad
(1) T3–T5 改为只凭位置的有限步形式：中点位移上链的环量与通量严格等于 $2A\im p$ 与 $\Delta A$，位置形式的 Kelvin 定理严格成立；(2) 每步先做 Kabsch 对齐，使 T3 的读数客观；(3) T8 的命题更正：方差单调衰减无条件成立，需要条件的是离散极值原理，而 Delaunay 剖分在电导 $\rho\kappa$ 变化时并不保证该条件；(4) 激波–等熵涡中旧的梯度能量占比（75\%–81\%）约三分之一来自平移，对齐后为 53\%；(5) 旧报告中已被推翻的掺混归因数值，统一以界面拉伸重推报告为准。

\textbf{\sffamily 评估结论}\quad
T1、T2、T6、T8、T9 理论完备、验证充分，可以单独构成一篇论文的主体；T3、T4 理论完备，但物理读数中只有局部量与区域积分得到验证，适合作为论文中的一节或另成一篇方法论文；T5、T7 不宜作为主要贡献。详见总结部分。
\vspace{0.3em}\end{minipage}}




%=====================================================================
\section{框架}

\paragraph{输入。}示踪点在快照上的位置 $x_i^n=x_i(t_n)$，$n=0,\dots,S-1$，以及一个参考时刻的密度 $\rho_i(t_c)$（只有 T2 需要）。速度、涡量与加速度不作为输入，只在解释结果时作为小步极限出现。

\paragraph{轨迹网络。}节点为示踪点；几何激活 $K^n$ 是第 $n$ 帧只从外边界剥离的 Delaunay 复形；动力学激活是窗口两端的互 kNN 图（AFRC 用）。关联矩阵 $B_1,B_2$ 即上边缘算子 $d_0,d_1$，满足 $d_1d_0=0$。只从外边界剥离，是为了不在内部开洞：只缺 3\% 的三角形，调和部分就会占到 68\%–84\% 的能量。

\paragraph{三点映射。}三角形 $t$ 从 $n$ 到 $n+1$ 的仿射映射写成 $F\zeta=p\zeta+q\bar\zeta$，
$p=\tfrac12[(F_{11}+F_{22})+i(F_{21}-F_{12})]$，$q=\tfrac12[(F_{11}-F_{22})+i(F_{21}+F_{12})]$。于是 $J=|p|^2-|q|^2$，极转角为 $\arg p$，$\mu=q/p$ 为 Beltrami 系数。两步合成为 $(p_2p_1+q_2\bar q_1,\ p_2q_1+q_2\bar p_1)$。三点之间的映射总是仿射的，所以这些量对任何位置数据都有定义。

\paragraph{原则。}(1) 计算中只出现相邻快照之间流映射的有限量；(2) 拉伸取点对，转动与环量取三角形，收支取物质区域；(3) 间断附近只报告积分量（区域环量、块映射、面积比），不报告逐三角形的密度量。

%=====================================================================
\section{T1 激波检测（C 通道）}

\paragraph{理论。}节点的单步星面积比 $\ln(A^\ast_{n+1}/A^\ast_n)=\int\theta\,dt$。光滑流中它是 $O(\Delta t)$；完整穿越密度比为 $\eta$ 的平面激波时为 $-\ln\eta=O(1)$，与 $\Delta t$ 无关。只要 $\sup|\theta|\Delta t\ll\tau<\ln\eta_{\min}$，就能把激波步与光滑压缩分开。带噪声时，比较台阶前后各两帧的平均对数星面积，阈值取 $\tau=\max(0.2,4\hat\sigma)$。

\begin{proposition}[可检出的最弱激波]\status{已证}
单步检出要求 $\ln\eta>\tau$。对完全气体 $\eta=(\gamma+1)M_s^2/[(\gamma-1)M_s^2+2]$，$\tau=0.2$、$\gamma=1.4$ 时要求 $M_s\gtrsim1.13$。更弱的激波，或一次穿越跨多个快照的粘性激波，需要把相邻步合并后再判断（第~\ref{pA:sec:block} 节的合并块）。
\end{proposition}

\paragraph{合并穿越块（补充）。}\label{pA:sec:block}三角形的热步（$|\ln J_t|>0.05$）用合成律并成块 $(p_B,q_B)$。完整穿越平面激波的块是单轴映射，满足 $\arg p_B=0$ 与 $|\mu_B|=\tanh(\tfrac12|\ln J_B|)$，与激波倾角无关。这条曲线可作为 C 通道的独立复核。

\paragraph{验证。}
\begin{center}\small
\begin{tabular}{lcccc}\toprule
& Sod & 爆炸波 & 激波–气泡 & 激波–涡\\\midrule
检出率 & 1.00 & 1.00 & 1.00 & 1.00（以 DNS 密度跃变 15\% 为准时 0.96）\\
误报率 & $10^{-4}$ & 0.016 & 0.006 & 0\\\bottomrule
\end{tabular}\end{center}
不可压算例中没有误报。合并块的签名在 Sod 中与单轴曲线一致到 $10^{-16}$，在爆炸波中中位偏离 0.011，在激波–涡中为 0.018（激波被涡弯曲）；逐步求和的伪转角 95 分位，在合并后由 0.092 降到 0.0057 rad（Sod）、由 0.29 降到 0.029 rad（爆炸波）。爆炸波的误报来自激波后仍在持续的压缩。

\begin{figure}[!htbp]\centering
\includegraphics[width=\linewidth]{figT1_shock.pdf}
\caption{(a) Sod 中每个示踪点被标记的热步数；(b)(c) 跨激波伪转角：逐步求和与合并块；(d) 三个激波算例的合并块签名，黑线为单轴曲线。}\label{pA:fig:t1}
\end{figure}

\paragraph{状态。}理论完备，验证充分。缺口：弱激波下限只做了理论估计；粘性激波跨多帧的情形只在 SBLI 中才会出现，尚未检验。

%=====================================================================
\section{T2 物质账本与质量测度}

\paragraph{理论。}示踪点质量 $m_i=\rho_i(t_c)A^\ast_i(t_c)/3$ 在窗口内守恒；密度 $\rho_i(t)=\rho_i(t_c)e^{-\ln J_i(t)}$，$\ln J_i=\sum_k\ln[A^\ast_{i,k}(t_{k+1})/A^\ast_{i,k}(t_k)]$ 为逐点锚定的星面积比之和。锚定星处于激波模板内时推迟重锚，使每段的两端分别位于激波前后，面积比跨激波望远相消。带噪声时有三处必要修正：锚定剖分取自相邻帧、面积在当前帧测量；星面积取有向面积；C 通道只认持续的台阶。

\begin{proposition}[离线追踪的误差恒等式]\status{已证+数值}
离线轨迹是“快照加线性时间插值”这一速度场的流映射。它与真实流映射之差满足 $\delta\ln J=\nabla\cdot\delta x$（一阶）。当快照激波 Courant 数 $C_s=U_s\Delta t_s/\Delta x$ 远大于激波数值厚度时，区间内的 $\delta\ln J\to\ln\eta$，与 $\Delta t_s$ 无关。适用条件为 $C_s\lesssim2$。
\end{proposition}

\paragraph{验证。}
\begin{center}\small
\resizebox{\linewidth}{!}{\begin{tabular}{lcccccc}\toprule
& 双环流 & Lamb–Oseen & Sod & 爆炸波 & 激波–气泡 & 激波–涡\\\midrule
密度中位误差 & 1.35\% & 0.48\% & $<10^{-14}$ & 0.58\% & 0.73\% & 0.36\%\\
95 分位 & 6\% & 18\% & 1\% & 35\% & 32\% & 4\%\\
0.1$h$ 噪声下的偏差 & +1.61\% & +0.67\% & +0.17\% & +0.91\% & +0.43\% & +0.43\%\\\bottomrule
\end{tabular}}\end{center}
95 分位出现在激波正扫过的星上、接触间断两侧，以及激波–气泡界面上。离线追踪在激波–涡中使密度中位误差从 0.33\% 升到 11.8\%；$\delta\ln J$ 与 $\nabla\cdot\delta x$ 的相关为 0.72。

\paragraph{状态。}理论完备，验证充分。缺口：(1) 需要一个参考时刻的密度，严格说不是“只凭位置”；(2) 激波–气泡界面上，账本压缩比 DNS 插值高 0.05–0.09，旧报告归因于 DNS 的数值扩散，但这只由 DNS 内部的空间分布推断，缺少更高分辨率 DNS 的直接确认；(3) 两份报告中同一量（$t=1.4$ 界面压缩）有 0.450 与 0.501 两个值，来自账本的噪声修订前后，成文时只能保留修订后的版本。

%=====================================================================
\section{T3 边信号的 Hodge 分解与局部读数}

\subsection{位置形式}
\begin{definition}[中点位移上链]
$d_e^{n+1/2}=\tfrac12(\Delta x_a+\Delta x_b)\cdot\tfrac12(q_e^n+q_e^{n+1})$，$\Delta x=x^{n+1}-x^n$，$q_e=x_b-x_a$。
\end{definition}
\begin{theorem}[有限步环量与通量]\status{已证+数值}
对任意三角形与任意一步，$(d_1d)_t=2A_t^n\im p_t$，法向形式的通量等于 $A_t^{n+1}-A_t^n$。
\end{theorem}
\begin{proof}
三点映射 $x^{n+1}=FX+c$。位移对 $X$ 线性，中点边向量为 $\tfrac12(I+F)Q$，梯形规则精确，由 Stokes 公式得 $A^n(K_{12}-K_{21})$，$K=(F-I)^\top\tfrac12(I+F)$，其反对称部分为 $\tfrac12(F^\top-F)$，故等于 $2A^n\im p$。通量：$A(s)=A^n\det(I+s(F-I))$ 是 $s$ 的二次式，$A'(s)$ 线性，$A^{n+1}-A^n=A'(\tfrac12)$，即中点构形上的通量。
\end{proof}
旧定义 $f_e=\tfrac12(u_a+u_b)\cdot(x_b-x_a)$ 是 $d_e/\Delta t$ 的小步极限；旧命题要求速度在三角形上线性，新定理不要求。由此得到只凭位置的步涡量 $\hat\omega_t=2A^n_t\im p_t/(\bar A_t\Delta t)$ 与步膨胀 $\hat\theta_t=\ln J_t/\Delta t$。

\subsection{分解、度量与客观性}
在内积 $\langle\cdot,\cdot\rangle_1$ 下 $C^1=\operatorname{im}d_0\oplus\operatorname{im}\delta_2\oplus\ker L_1$，$\dim\ker L_1=\beta_1(K)$。度量取 Galerkin（P1/Crouzeix–Raviart）：弱散度 $\mathrm{div}_i=-\tfrac{3}{A^\ast_i}\sum_{t\ni i}A_t\bar u_t\cdot\nabla\lambda_{i,t}$ 对任意剖分上的仿射场精确；对角 cotan 度量的梯度泄漏中位数为 24\%，Galerkin 为 $7.5\times10^{-4}$。质量只进账本，不进内积：以 $\rho$ 加权时余正合空间变为 $\nabla\cdot(\rho u)=0$，无散的变密度流会出现虚假的梯度部分。

\begin{proposition}[Kabsch 对齐]\status{已证+数值}
每步先求刚体运动 $(R^n,b^n)=\arg\min\sum_im_i|x_i^{n+1}-R^nx_i^n-b^n|^2$，再把 $x^{n+1}$ 拉回第 $n$ 帧参考系。由此计算的 $p,q,J,\mu,d_e$ 对任意 $x\mapsto Q(t)x+c(t)$ 不变。
\end{proposition}
均匀平移是纯梯度场，所以不对齐时梯度能量占比依赖参考系。对齐所用的转动是系综量，所以“客观”指对刚体参考系变换不变，并非局部。

\subsection{验证}
\begin{center}\small
\resizebox{\linewidth}{!}{\begin{tabular}{lcccccc}\toprule
& 双环流 & Lamb–Oseen & Sod & 爆炸波 & 激波–气泡 & 激波–涡\\\midrule
涡量相关（旧，速度上链） & 1.000 & 0.998 & 伪 0.007 & 伪 0.000 & 0.996 & 0.999\\
$\hat\omega$ 相关（新，位置形式） & 0.97–0.999 & 0.93–0.9995 & — & — & 未重算 & 0.976–0.9999$^a$\\
散度相关（中间 / 末时刻） & 伪 0.011 & 伪 0.016 & 0.99 / 0.86 & 0.89 / 0.56 & 0.77 / 0.90 & 0.87 / 0.80\\
梯度能量（原始） & 0.01 & 0.10 & 1.00 & 1.00 & 0.89 & 0.75–0.81\\
梯度能量（Kabsch 对齐） & 未重算 & 0.037$^b$ & 未重算 & 未重算 & 未重算 & 0.53\\
0.1$h$ 噪声下涡量相关 & 0.93 & 0.89 & — & — & 0.67 & 0.78\\\bottomrule
\end{tabular}}\\[0.2em]{\footnotesize $^a$ 去掉跨激波三角形；包含时最低 0.47。\quad $^b$ 单帧，原始为 0.030。}\end{center}

\begin{figure}[!htbp]\centering
\includegraphics[width=\linewidth]{figT2_sv.pdf}
\caption{激波–等熵涡（DNS）。(a) 步涡量与 DNS，绿色为跨激波三角形；(b) 涡核物质区域的环量，绿带为激波位于区域内的步；(c) 位移上链 Hodge 分解的能量占比：原始计算加一个 0.01 的平移就从 75\% 变为 94\%，Kabsch 对齐后为 53\% 且不变。}\label{pA:fig:t3}
\end{figure}

\paragraph{状态。}分解定理本身是标准的有限维 Hodge 定理，位置形式与对齐是本次补全。得到验证的是局部读数（涡量、散度）与区域积分。三个问题需要在成文时正视：
\begin{enumerate}
\item \textbf{能量占比不能作为物理读数。}即使对齐，梯度/余正合的划分也是非局部的：鞍点区局部无旋，却在 Biot–Savart 意义下属于余正合部分；自由边界的法向通量会被归入梯度部分（Lamb–Oseen 中的 4\%–10\%）。旧报告第四部分把能量占比解释为“由物理决定的差别”，这一说法应删除。
\item \textbf{激波附近的逐三角形涡量是伪量。}跨激波三角形使相关降到 0.47（激波–涡），这些三角形上只能报告积分量。
\item \textbf{散度明显弱于涡量}（0.56–0.99），在压缩已经衰减的时刻尤其弱。
\end{enumerate}

%=====================================================================
\section{T4 离散 Kelvin 定理与环量收支}

\begin{theorem}[位置形式的离散 Kelvin 定理]\status{已证+数值}
记 $\Gamma_t^{n+1/2}=(d_1d^{n+1/2})_t$，二阶位移上链 $g_e^n=\tfrac12(\Delta^2x_a+\Delta^2x_b)\cdot\tilde q_e$，$\tilde q_e=q_e^n+c(q_e^{n+1}-2q_e^n+q_e^{n-1})$，$c\in\mathbb R$。则对任意三帧位置严格成立 $\Gamma_t^{n+1/2}-\Gamma_t^{n-1/2}=(d_1g^n)_t$。
\end{theorem}
\begin{proof}
记 $u^\pm$ 为前后半步的边平均位移，$\delta^\pm$ 为前后半步的边向量变化，则一条边上的差为 $(u^+-u^-)\cdot q^n+\tfrac12(u^+\!\cdot\delta^++u^-\!\cdot\delta^-)$。后一项等于 $\tfrac14(|\Delta x_b^+|^2-|\Delta x_a^+|^2+|\Delta x_b^-|^2-|\Delta x_a^-|^2)$，是恰当差，绕三角形求和为零；$c\neq0$ 时多出的项 $\tfrac c2(\Delta^2x_a+\Delta^2x_b)\cdot(\Delta^2x_b-\Delta^2x_a)$ 同样是恰当差。
\end{proof}
连续极限下 $(d_1g)_t/\Delta t^2\to\oint_{\partial t}a\cdot dx=\int_t\rho^{-2}\nabla\rho\times\nabla p\,dA+$ 黏性项。物质区域上的和只依赖边界上的上链，与内部剖分无关，所以每隔若干帧重新剖分内部不改变区域环量。激波步上 $\Delta^2x$ 是有限的速度跳跃乘 $\Delta t$，“冲击部分”有严格定义；在合并块上计算可以避免逐步求和的伪转动。

\paragraph{验证。}恒等式在全部算例（含 DNS）上的残差不超过 $5\times10^{-14}$，且不随 $\Delta t$ 变化。区域环量与参照（同一组三角形上积分参照涡量）的误差：双环流 2.5\%，Lamb–Oseen 0.5\%，激波–气泡 1.3\%（全程最大 6.1\%），激波–涡 0.0\%–0.8\%；无旋算例的伪环量为 $3\times10^{-4}$（以 $u_{\max}|\partial R|$ 归一）。激波–气泡中，DNS 重气体的欧拉环量是内部物质轮廓的 3 倍，说明斜压沉积的环量主要在界面涡片上。

\paragraph{状态。}理论完备且严格。缺口：(1) 区域收支中的“冲击部分 / 光滑部分”拆分没有与 DNS 的斜压力矩 $\int\rho^{-2}\nabla\rho\times\nabla p\,dA$ 直接比较（统一报告没有，早期 WP3 有力矩去噪的对比，成文时需要并入）；(2) Sod 中“冲击部分 / 总变化 = 471\%”的分母接近零，应从表中删去；(3) 拆分应改在合并块上定义，目前按“激波位于区域内的步”划分，会把这些步内的光滑沉积也算进冲击部分。

%=====================================================================
\section{T5 Dirac 结构}

\paragraph{理论。}$D=\begin{psmallmatrix}0&\delta_1&0\\d_0&0&\delta_2\\0&d_1&0\end{psmallmatrix}$，$D^2=\mathrm{diag}(L_0,L_1,L_2)$，$D$ 关于物理度量自伴；节点—梯度与余正合—三角形的谱配对给出 Hodge 分解；AFRC 的结构项 $\mathrm{struct}(e)$ 是组合权下 $L_1$ 的对角占优余量。块映射上的有限残差为 $r_B=\ln J_B-(\operatorname{tr}F_B-2)$：光滑流中为 $O(\Delta t^2)$，平面激波中为 $1-\eta^{-1}-\ln\eta$。

\paragraph{验证。}旧残差 $\theta_C-\mathrm{div}$ 在光滑区与激波处之比：Sod 0.12 对 3.85，爆炸波 0 对 0.24，激波–气泡 0.002 对 0.005，激波–涡 0 对 0.014。DNS 中对比度弱。

\paragraph{状态。}代数结构完备，物理检验薄弱。有限残差与 T1 的签名 $(\ln J,|\mu|)$ 冗余；声学 Dirac 方程与 DESP 联合滤波没有定量检验。建议在论文中作为附录或备注（说明 AFRC 与 LHD 共用同一算子族），不作为结果。

%=====================================================================
\section{T6–T7 物质输运与转移算子}

\paragraph{理论。}控制体质量 $\hat M_\Omega(t)=\sum_im_i\chi_\Omega(x_i(t))$；单向输运 $F^\to=\sum_im_i[x_i(t_0)\notin\Omega][x_i(t_1)\in\Omega]$，误差 $O(h/L)$。按质量加权的 Ulam 矩阵 $P_{ij}\approx\sum_km_k[x_k(t_0)\in B_i][x_k(t_1)\in C_j]$，行和、列和都精确；按面积加权时列和是“被推前的初始面积”，在可压缩流中与 $C_j$ 的实际面积相差 $\int(J-1)\,dA$。组分必须按质量分数计算。

\paragraph{验证。}截面通量误差：双环流 0.33\%，Lamb–Oseen 0.30\%，Sod 0.91\%，爆炸波 0.16\%，激波–气泡 1.37\%；双环流的单向交换随分辨率一阶收敛（$-2.1\%$、$-0.8\%$、$-0.34\%$）。卷起涡内的组分按计数偏差 2.5 倍，按质量偏差约 23\%（主要来自 DNS 被动标量的数值扩散）。转移算子：Sod、爆炸波、激波–涡中 $\sigma_2=1$（没有混合，退化）；Lamb–Oseen 主导相干集与涡核的 Jaccard 只有 0.11；按质量与按面积加权的 $\sigma_2$ 只差 0.0006（激波–气泡）。

\paragraph{状态。}T6 理论完备、验证充分，但估计量本身简单，贡献在于“测度取质量”。T7 定义完备，但六个算例中三个退化、一个给出错误的主导集，按质量加权的效果也很小，不宜作为主要贡献。

%=====================================================================
\section{T8 网络上的拉格朗日对流—扩散}

\paragraph{理论。}由 $\rho\,Dc/Dt=\nabla\cdot(\rho\kappa\nabla c)$ 的物质弱形式，P1 元并集中质量：
\begin{equation}
M\dot c=-K(t)c,\qquad M=\mathrm{diag}(m_i),\qquad K_{ij}=\sum_{T\ni i,j}(\rho\kappa)_TA_T\nabla\lambda_i\cdot\nabla\lambda_j .
\end{equation}
质量矩阵与时间无关；对流由示踪点精确完成。

\begin{proposition}[守恒、能量衰减与极值原理]\status{已证}\label{pA:prop:t8}
(i) $\tfrac{d}{dt}\sum_im_ic_i=0$。(ii) 按质量加权的方差 $V=\sum_im_i(c_i-\bar c)^2$ 满足 $\dot V=-2\sum_T(\rho\kappa)_TA_T|\nabla c_T|^2\le0$，\textbf{无条件}成立；隐式 Euler $(M+\Delta tK^n)c^{n+1}=Mc^n$ 下 $V^{n+1}\le V^n$ 同样无条件成立，每个快照换剖分也不影响。(iii) 离散极值原理（$c$ 不产生新的极值）要求 $K_{ij}\le0$（$i\ne j$）。由 $K_{ij}=-\tfrac12\sum_{T\ni ij}(\rho\kappa)_T\cot\theta_T$（$\theta_T$ 为 $T$ 中对着边 $ij$ 的角），Delaunay 只保证 $\cot\theta_1+\cot\theta_2\ge0$；当 $(\rho\kappa)$ 在两个三角形上不同、且其中一个角为钝角时，$K_{ij}$ 可以为正。充分条件是：若 $\theta_2>\pi/2$，则 $(\rho\kappa)_1\ge(\rho\kappa)_2$；外边界边要求对角不为钝角。
\end{proposition}
\begin{proof}
(i) $K$ 对称、行和为零。(ii) $c^\top Kc=\sum_T(\rho\kappa)_TA_T|\nabla c_T|^2\ge0$；隐式步两边乘 $c^{n+1}$ 并用 Cauchy–Schwarz。(iii) 直接计算。
\end{proof}
旧版命题把方差衰减写成以 $K_{ij}\le0$ 为条件，并称 Delaunay 内部边满足该条件；两处都应按命题~\ref{pA:prop:t8} 更正。两份旧文档中 $K_{ij}$ 符号约定不一致的问题也由此统一。若需要严格的极值原理，可把电导取在边上（例如两侧三角形的调和平均），此时 $K_{ij}=-\tfrac12(\rho\kappa)_e(\cot\theta_1+\cot\theta_2)\le0$ 在 Delaunay 内部边上成立；这一改法尚未做数值检验。

\paragraph{体积形式的不一致。}若用 $M=\mathrm{diag}(A^\ast_i/3)$ 与电导 $\kappa$，离散的是 $Dc/Dt=\kappa\nabla^2c$，比正确方程少 $\kappa\nabla\ln\rho\cdot\nabla c$，且守恒的是 $\int c\,dA$。不可压时两者等价。

\paragraph{验证。}末时刻相对 $L^2$ 误差（质量形式 / 体积形式 / 纯对流）：双环流 3.2\% / 2.8\% / 66\%，Lamb–Oseen 1.6\% / 1.3\% / 35\%，Sod 0.2\% / 1.8\% / 7\%，爆炸波 2.3\% / 22.3\% / 109\%，激波–气泡 5.1\% / 25.2\% / 55\%。可压缩算例中体积形式的误差大 5–10 倍；方差比与参照相差在 1\% 以内；激波–气泡中标量质量严格守恒。

\paragraph{状态。}理论完备（本稿更正了命题的条件），验证充分，可压缩与不可压的对比清楚。缺口：极值原理在激波附近是否被破坏、破坏多少，尚未统计。

%=====================================================================
\section{T9 掺混：拉伸与压缩}

\begin{proposition}[梯度增长的精确分解]\status{已证}
二维 Nanson 公式 $|FT|=J|F^{-\top}N|$ 给出，被动标量（$\kappa\to0$，$c_0=g(\phi_0)$）沿每个物质点满足
$\ln(|\nabla c|/|\nabla c_0|)=\ln\lambda-\ln J$，$\lambda=|FT|$。这是代数恒等式，不需要“薄界面元”假设。
\end{proposition}

\paragraph{估计量。}(1) 点对估计的偏差只依赖方向，与示踪点间距 $h$ 无关，加密不收敛；按点取中位数只在 $\dot\gamma t<2\cot\alpha_{\max}\approx4.8$ 内无偏。(2) 链式张量估计：近邻上二次最小二乘拟合增量变形梯度，按账本规则串联，是一致估计，误差被各向异性 $A=\sigma_{\max}(F)/\lambda$ 放大。(3) Nanson 余向量估计：$\hat\lambda=\hat J|\widehat{\nabla_x\phi_0}|/|\nabla_X\phi_0|$，不做时间积分，但在物质层薄于 $h$ 时偏低。当界面既是强剪切面、物质层又薄于 $h$ 时，两者都会失效（充分条件；数值见第六部分），这是只凭位置的信息极限。(4) 涡片上的拉伸本来就是两值的：$\sigma^+-\sigma^-=\partial_s\Delta u_t$。

\paragraph{验证。}
\begin{center}\small
\resizebox{\linewidth}{!}{\begin{tabular}{lcccccc}\toprule
& 双环流 & Lamb–Oseen & Sod & 爆炸波 & 激波–气泡 & 激波–涡\\\midrule
拉伸（链式张量 / 参照） & 1.708 / 1.711 & $-0.551$ / $-0.310$ & 0.000 / 0 & 0.576 / 0.577 & 0.285 / 0.390$^a$ & $-0.231$ / $-0.224$\\
压缩（账本 / 参照） & 0.000 / 0 & $-0.005$ / 0 & 0.754 / 0.754 & 0.637 / 0.638 & 0.45–0.50 / 0.36$^b$ & 0.396 / 0.398\\\bottomrule
\end{tabular}}\\[0.2em]{\footnotesize $^a$ 涡片两侧参照曲线的范围为 [0.06, 0.39]。\quad $^b$ DNS 在界面附近有数值扩散，不能作为压缩参照（见 T2 状态）。}\end{center}
链式张量的逐点误差随分辨率收敛：Lamb–Oseen 为 0.78、0.26、0.06（二阶），爆炸波已接近机器精度；点对法停在 0.8–1.0 不收敛。

\begin{figure}[!htbp]\centering
\includegraphics[width=0.92\linewidth]{fig_sb_attr.png}
\caption{激波–重气泡界面的掺混归因（取自界面拉伸重推报告）。(a) 压缩项、拉伸项与总增长的时间历程；(b) 新旧归因对比：界面在压缩后先被切向缩短，到 $t=1.4$ 压缩仍占总增长的 54\%–62\%。}\label{pA:fig:t9}
\end{figure}

\paragraph{物理结论。}激波–重气泡中，激波把气泡沿流向压扁，与激波大致平行的界面段被切向缩短，拉伸项在 $t\approx0.7$ 达到 $-0.4$ 左右，抵消了大部分压缩带来的梯度增长；卷起之后拉伸项才转正。到 $t=1.4$ 压缩仍是主要贡献，界面长度比为 2.2 倍。旧报告“压缩只占约 30\%、卷起的拉伸逐渐接管、长度 5.3 倍”的结论已被推翻。

\paragraph{状态。}理论完备，验证充分，并给出了一个修正旧结论的物理结果。缺口：Lamb–Oseen 圆环（$A$ 中位数 13）仍在信息极限附近；激波算例中账本对位置噪声敏感，统一流水线的修正在压缩项上的效果需要单列。

%=====================================================================
\section{T10 位置噪声}

每帧加独立高斯噪声 0.05$h$ 与 0.1$h$，重跑 T1–T4 与 T9。0.1$h$ 时密度偏差不超过 1.6\%；最敏感的是逐点涡量（DNS 算例中相关降到 0.67–0.78）；区域环量、拉伸与压缩基本稳定。噪声来源的分析表明：位置噪声在固定物质三角形上近似望远相消，逐帧重剖时按 $\sqrt n$ 累积。

%=====================================================================

\clearpage\part{第二部分　框架重组与补充验证}

\noindent\fcolorbox{accent!40}{tint}{\begin{minipage}{0.965\linewidth}\vspace{0.3em}
\textbf{\sffamily 本稿做了什么}\quad 按上一份评估列出的缺口，补做六项工作：(1) 两个 DNS 各算三档网格（$\Delta x=0.01,0.005,0.0025$，示踪点不变），确认界面压缩的参照与密度 95 分位误差的来源；(2) 密度与 DTFE、SPH 对比，扩散与 SPH 传导对比；(3) 账本与链式张量的参数敏感性；(4) 稀疏示踪点、断轨与实际量级噪声下的检验；(5) 新增单模 Richtmyer–Meshkov（RM）DNS 作为复杂应用；(6) 统一全部文档中的数值。另外回答“其余 T 是否融入、没有 T3–T5 还能否叫 LHD”。

\textbf{\sffamily 主要结论}\quad
\textbf{框架}：论文 A 的成熟部分是同一个 Dirac 算子在质量度量下的几个块，可以保留 LHD 的名字（第~1 节）。
\textbf{DNS 加密}：激波–等熵涡的 95 分位来自账本在涡核内的误差，与 DNS 无关；激波–重气泡卷起之前的界面差异来自 DNS 数值扩散，卷起之后（$t=1.4$）界面带上两者都不收敛，旧的“压缩占 54\%–62\%”需降级为待复核（第~2 节）。
\textbf{基线}：只凭位置的密度比 SPH 准约一个数量级、比 DTFE 准 1–2 个数量级；扩散在三个算例上比 SPH 准 2–5 倍。
\textbf{参数}：只有“最多推迟重锚次数”敏感，不推迟时爆炸波误差增大约 18 倍。
\textbf{实际数据}：0.02$h$ 噪声下中位误差约 1.5\%–2\%；稀疏、断轨下仍可用。
\textbf{RM}：振幅与 DNS 相差 $\le$1.2\%，增长率与 Richtmyer 模型相差 3\%，模型参数全部取自示踪点位置。
\vspace{0.3em}\end{minipage}}



\section{框架重组：质量度量下的 Hodge–Dirac}
\subsection{三个命题}
记第 $n$ 帧的剥离 Delaunay 复形为 $K^n$，上边缘 $d_0,d_1$，度量 $M_0,M_1,M_2$，$\delta_1=M_0^{-1}d_0^\top M_1$。三点映射写成 $F\zeta=p\zeta+q\bar\zeta$。

\begin{proposition}[复上链恒等式]\status{已证+数值}\label{pB:prop:cplx}
对边 $e=(a\to b)$ 与步 $n\to n+1$，定义复 1‑上链 $c_e=\overline{\bar u_e}\;\bar q_e$，$\bar u_e=\tfrac12(\Delta x_a+\Delta x_b)$，$\bar q_e=\tfrac12(q_e^n+q_e^{n+1})$（二维向量与复数等同）。则对任意三角形与任意一步严格成立
\begin{equation}
(d_1c)_t=2A_t^n\im p_t+i\,(A_t^{n+1}-A_t^n).
\end{equation}
\end{proposition}
\begin{proof}
$\re(\bar u\,q)=u\cdot q$，$\im(\bar u\,q)=u_xq_y-u_yq_x=u\cdot R_{-90^\circ}q$。实部即中点位移上链的环量，虚部即法向通量；两者分别等于 $2A^n\im p$ 与 $A^{n+1}-A^n$（T3 定理，证明见 T1–T9 定稿）。
\end{proof}
于是 \textbf{T4}（Kelvin 环量）是 $\re(d_1c)$，\textbf{T2}（账本）的单步增量是 $\ln(1+\sum_{t\ni i}\im(d_1c)_t/A^\ast_i)$：二者是同一个有限步量的两个分量，不需要速度，也不依赖步长。数值上两部分分别验证到 $3\times10^{-14}$ 与 $2\times10^{-14}$。

\begin{proposition}[扩散算子是 $D^2$ 的节点块]\status{已证}\label{pB:prop:L0}
取 $M_0=\mathrm{diag}(m_i)$（示踪点质量），$M_1$ 为以 $(\rho\kappa)_T$ 加权的 P1 Galerkin 度量，则 T8 的半离散方程 $M\dot c=-Kc$ 就是 $\dot c=-L_0c$，$L_0=\delta_1d_0=M_0^{-1}d_0^\top M_1d_0$。体积形式对应 $M_0=\mathrm{diag}(A^\ast_i/3)$、$M_1$ 以 $\kappa$ 加权。
\end{proposition}
所以“质量形式比体积形式准 5–10 倍”（可压缩算例）是对度量选择的直接检验：物理测度决定 Hodge–Dirac 的度量，而不是数值细节。

\begin{proposition}[其余组成的位置]\status{已证}
T9 的 Nanson 余向量估计是 $d_0$ 作用在初始标签 0‑上链 $\phi_0$ 上；链式张量估计是三点映射合成律在近邻上的最小二乘推广；T1 的激波签名是合并块满足 $|\mu_B|=\tanh(\tfrac12|\ln J_B|)$、$\arg p_B=0$；T6 是质量 0‑链的推前。
\end{proposition}

\subsection{各组成部分的去留}
\begin{table}[!htbp]\centering\small
\caption{组成部分与 Dirac 算子的对应，以及在论文 A 中的位置。}
\begin{tabularx}{\linewidth}{@{}l L l@{}}\toprule
组成 & 在 $D$（质量度量）中的位置 & 去向\\\midrule
T1 激波检测 & 节点上的面积跳跃；块映射的签名 & 论文 A\\
T2 物质账本 & $\im(d_1c)$ 在星上求和 & 论文 A（核心）\\
T3 局部读数 & $\re(d_1c)/(\bar A\Delta t)$，$\im(d_1c)/(\bar A\Delta t)$ & 论文 A（只保留局部量）\\
T3 全局分解 & $D$ 的谱分割、能量占比 & 另一篇（拓扑信号）\\
T4 Kelvin & $\re(d_1c)$ 与二阶上链 $d_1g$ & 论文 A\\
T5 Dirac 谱、声学、DESP & $D$ 的谱与联合滤波 & 另一篇（拓扑信号）\\
T6 质量输运 & 质量 0‑链的推前 & 论文 A\\
T7 转移算子 & 推前的 Ulam 矩阵 & 论文 A 一段或删去\\
T8 扩散 & $L_0=\delta_1d_0$，度量取质量与 $\rho\kappa$ & 论文 A（核心）\\
T9 掺混 & Nanson：$d_0\phi_0$ 与账本；链式 $(p,q)$ & 论文 A（核心）\\
脊分类 & 边时间序列与账本 & 第三项工作\\\bottomrule
\end{tabularx}\end{table}

\subsection{论文 A 的创新点}
\begin{enumerate}
\item \textbf{有限步严格。}环量、面积变化、Kelvin 收支都是位置的代数恒等式，不需要速度、不依赖步长，跨激波时用合并块仍然成立。已有的拉格朗日轨迹网络方法建立在图（0、1 维）上，没有 2‑胞腔上的这类恒等式。
\item \textbf{度量即测度。}可压缩流中质量度量与面积度量给出可区分的结果：扩散误差相差 4–10 倍，组分按计数偏差 2.5 倍。
\item \textbf{精度。}只凭位置的密度比 DTFE、SPH 准 1–2 个数量级（第~\ref{pB:sec:base} 节）。
\item \textbf{物理结果。}激波–重气泡界面在激波扫过后先被切向缩短、抵消大部分压缩（$t\le0.6$，对 DNS 网格稳健）；单模 RM 到 $ka\approx0.6$ 时界面梯度增长约 94\% 来自压缩，振幅增长率与 Richtmyer 模型相差 3\%，模型参数全部取自示踪点位置（第~\ref{pB:sec:rm} 节）。
\end{enumerate}
\section{DNS 加密：界面压缩的参照与密度 95 分位的来源}\label{pB:sec:dns}
两个 DNS 各算三档网格 $\Delta x=0.01,0.005,0.0025$，示踪点初始位置与间距（$h=0.005$）不变，在求解器内用同样的 Heun/双线性格式推进。比较同一组示踪点上的账本压缩 $-\ln J$ 与 DNS 插值密度比 $\ln(\rho/\rho_0)$，按初始到界面的有符号距离 $\phi_0$ 分带（负值为重气体一侧）。
\begin{table}[!htbp]\centering\small\caption{激波–重气泡，$t=0.6$：各带的账本压缩 / DNS 插值压缩。}
\begin{tabular}{lccc}\toprule
$\phi_0$ 区间 & $\Delta x=0.01$ & $\Delta x=0.005$ & $\Delta x=0.0025$\\\midrule
$[-0.06,-0.03]$ & 0.594 / 0.594 & 0.600 / 0.611 & 0.603 / 0.609\\
$[-0.03,-0.015]$ & 0.538 / 0.463 & 0.548 / 0.528 & 0.552 / 0.555\\
$[-0.015,-0.005]$ & 0.506 / 0.371 & 0.514 / 0.387 & 0.506 / 0.463\\
$[-0.003,0.003]$ & 0.491 / 0.436 & 0.481 / 0.426 & 0.300 / 0.436\\
$[0.005,0.015]$ & 0.495 / 0.597 & 0.476 / 0.607 & 0.409 / 0.549\\
$[0.015,0.03]$ & 0.486 / 0.572 & 0.468 / 0.522 & 0.464 / 0.489\\
$[0.03,0.06]$ & 0.490 / 0.530 & 0.487 / 0.503 & 0.485 / 0.489\\
\bottomrule\end{tabular}\end{table}
\begin{table}[!htbp]\centering\small\caption{激波–重气泡，$t=1.4$：各带的账本压缩 / DNS 插值压缩。}
\begin{tabular}{lccc}\toprule
$\phi_0$ 区间 & $\Delta x=0.01$ & $\Delta x=0.005$ & $\Delta x=0.0025$\\\midrule
$[-0.06,-0.03]$ & 0.578 / 0.505 & 0.585 / 0.568 & 0.581 / 0.581\\
$[-0.03,-0.015]$ & 0.538 / 0.340 & 0.559 / 0.441 & 0.559 / 0.517\\
$[-0.015,-0.005]$ & 0.485 / 0.218 & 0.527 / 0.210 & 0.516 / 0.298\\
$[-0.003,0.003]$ & 0.457 / 0.416 & 0.468 / 0.363 & -0.109 / 0.358\\
$[0.005,0.015]$ & 0.536 / 0.813 & 0.533 / 0.956 & 0.426 / 0.854\\
$[0.015,0.03]$ & 0.582 / 0.835 & 0.578 / 0.717 & 0.511 / 0.655\\
$[0.03,0.06]$ & 0.578 / 0.649 & 0.556 / 0.620 & 0.527 / 0.600\\
\bottomrule\end{tabular}\end{table}
\begin{table}[!htbp]\centering\small\caption{激波–重气泡：密度误差随 DNS 网格的变化。}
\begin{tabular}{lccc}\toprule
量 & $\Delta x=0.01$ & $\Delta x=0.005$ & $\Delta x=0.0025$\\\midrule
$t=0.6$ 中位误差 & 0.51\% & 0.28\% & 0.19\%\\
$t=0.6$ 95 分位 & 13.5\% & 10.4\% & 6.1\%\\
$t=0.6$ 95 分位（离界面 $>0.03$） & 5.1\% & 2.0\% & 1.0\%\\
$t=1.4$ 中位误差 & 1.3\% & 0.67\% & 0.45\%\\
$t=1.4$ 95 分位 & 29.8\% & 29.7\% & 31.3\%\\
$t=1.4$ 95 分位（离界面 $>0.03$） & 16.7\% & 9.9\% & 7.1\%\\
$x>1$ 的重气体质量：网络 / DNS & 0.2864 / 0.2889 & 0.2888 / 0.2916 & 0.2897 / 0.2931\\\bottomrule\end{tabular}\end{table}
\begin{table}[!htbp]\centering\small\caption{激波–等熵涡：密度误差随 DNS 网格的变化（$t=0.6$）。}
\begin{tabular}{lccc}\toprule
量 & $\Delta x=0.01$ & $\Delta x=0.005$ & $\Delta x=0.0025$\\\midrule
中位误差 & 0.54\% & 0.34\% & 0.27\%\\
95 分位 & 3.9\% & 3.4\% & 3.5\%\\
涡核内中位 / 95 分位 & 1.8\% / 7.6\% & 1.6\% / 7.0\% & 1.5\% / 7.4\%\\
误差最大的 5\% 中位于涡核的比例 & 96\% & 96\% & 93\%\\\bottomrule
\end{tabular}\end{table}
\paragraph{激波–等熵涡：95 分位来自账本在涡核内的误差。}DNS 网格从 0.01 加密到 0.0025，密度中位误差由 0.54\% 降到 0.27\%，说明中位误差中有一部分来自 DNS 参照本身；95 分位却保持在 3.4\%–3.9\%，误差最大的 5\% 中有 93\%–96\% 位于涡核（$r_0<2R_v$）内，涡核内的 95 分位为 7.4\%–7.6\%，不随 DNS 加密变化。所以本例的 95 分位是账本自身的误差：涡核内强差动旋转使锚定星在一个锚定周期内发生非仿射变形，星面积比不再是物质面积比。锚定周期取 3、5、8 帧时激波–涡的 95 分位为 2.9\%、3.4\%、4.7\%，印证了这一点。
\paragraph{激波–重气泡，卷起之前（$t=0.6$）：差异来自 DNS。}界面两侧 $0.005<|\phi_0|<0.03$ 的四个带上，账本与 DNS 插值压缩之差的平均值随 DNS 网格（0.01 / 0.005 / 0.0025）为 0.099 / 0.083 / 0.053；密度误差的 95 分位为 13.5\% / 10.4\% / 6.1\%，离界面 $>0.03$ 的 95 分位为 5.1\% / 2.0\% / 1.0\%，中位误差为 0.51\% / 0.28\% / 0.19\%。账本值在三档 DNS 轨迹上基本不变，DNS 值逐档向账本靠拢，两侧方向相反（重气体一侧 DNS 偏低、轻气体一侧偏高），这是接触间断被数值扩散抹宽、两种气体密度相互混合的特征。所以在卷起之前，账本与 DNS 的差异应归于 DNS 参照，账本给出的是示踪点所代表物质本身的压缩。
\paragraph{激波–重气泡，卷起之后（$t=1.4$）：离界面的带收敛，界面带不收敛。}重气体一侧 $[-0.06,-0.03]$ 与 $[-0.03,-0.015]$ 带上，DNS 压缩随网格为 0.505 / 0.568 / 0.581 与 0.340 / 0.441 / 0.517，逐档靠近账本（0.578 / 0.585 / 0.581 与 0.538 / 0.559 / 0.559）；离界面 $>0.03$ 的密度误差 95 分位由 16.7\% 降到 7.1\%，中位误差由 1.3\% 降到 0.45\%。但在 $|\phi_0|<0.015$ 内两者都不收敛：界面带 $[-0.003,0.003]$ 上账本压缩为 0.457 / 0.468 / -0.109，DNS 为 0.416 / 0.363 / 0.358；95 分位停在 29.8\% / 29.7\% / 31.3\%。原因在于这里的 DNS 是无粘的：涡片上的 Kelvin–Helmholtz 卷起随网格加密不断出现更细的结构，没有网格收敛的极限；$\Delta x=0.0025$ 时这些结构已经细于示踪点间距 $h=0.005$，锚定星内的映射不再是仿射的，账本也随之失效。这正是 T9 所说的信息极限（物质层薄于 $h$）。
\paragraph{结论与对旧结论的影响。}(1) 卷起之前，界面附近的账本—DNS 差异来自 DNS 的数值扩散，账本可信。(2) 卷起之后，界面带上的压缩参照无法由无粘 DNS 的加密确定，账本在界面带上也达到信息极限；离界面 0.015 以外两者收敛。(3) 因此旧报告中“$t=1.4$ 时压缩占总增长 54\%–62\%”依赖界面带上的压缩值，应降级为“依赖分辨率、需在粘性 DNS 上复核”；“激波扫过后界面先被切向缩短、抵消大部分压缩”（$t\le0.6$）不受影响。(4) 积分量对此不敏感：$x>1$ 的重气体质量（网络 / DNS）三档为 0.2864 / 0.2889；0.2888 / 0.2916；0.2897 / 0.2931，误差稳定在约 1\%。
\begin{figure}[!htbp]\centering\includegraphics[width=\linewidth]{figR_dns.pdf}\caption{(a)(b) 激波–重气泡各带的账本压缩（实线）与 DNS 插值压缩（虚线），颜色由浅到深为 DNS 网格由粗到细；(c) 两个 DNS 算例的密度误差随 DNS 网格的变化。}\end{figure}

\section{基线对比}\label{pB:sec:base}
\paragraph{密度。}三种方法用同一组示踪点质量 $m_i=\rho_0A^\ast_i(t_0)/3$：账本 $\rho_i=\rho_{0,i}e^{-\ln J_i}$（C 通道与推迟重锚，只凭位置）；DTFE（Schaap \& van de Weygaert 2000，$\rho_i=3m_i/A^\ast_i$）；自适应核 SPH 求和（三次样条，$k=24$）。
\begin{table}[!htbp]\centering\small\caption{末时刻密度的中位相对误差 / 95 分位（同一内部掩码）。}\begin{tabular}{lccc}\toprule 算例 & 账本 & SPH & DTFE\\\midrule
双环流 & 1.3\% / 5.7\% & 11.4\% / 65.3\% & 32.1\% / 254.3\%\\
Lamb–Oseen & 0.41\% / 13.0\% & 6.4\% / 49.2\% & 26.3\% / 167.6\%\\
Sod & $<10^{-6}$ / 0.97\% & 1.9\% / 7.1\% & 0.20\% / 37.8\%\\
爆炸波 & 0.58\% / 32.2\% & 6.8\% / 38.9\% & 20.2\% / 107.2\%\\
激波–涡 & 0.34\% / 3.4\% & 3.4\% / 32.2\% & 16.8\% / 97.7\%\\
激波–气泡 & 0.67\% / 29.7\% & 7.1\% / 43.6\% & 25.9\% / 141.7\%\\
RM & 0.73\% / 35.3\% & 4.6\% / 40.6\% & 17.1\% / 72.8\%\\
\bottomrule\end{tabular}\end{table}
在激波与接触间断附近（以精确解的局部跳跃定义），三者的 95 分位分别为：Sod 中账本 24.1\%、DTFE 33.3\%、SPH 42.6\%；爆炸波中 38.5\%、116.5\%、41.4\%。账本的中位误差在全部七个算例上比 SPH 小约一个数量级、比 DTFE 小 1–2 个数量级。DTFE 的误差主要来自 Delaunay 星在大变形下的重连：质量按初始星分配，当前星却已换了邻居；账本逐点锚定、跨激波望远相消，正是为此设计的。

\paragraph{扩散。}对比方法为 SPH 传导（Brookshaw / Cleary–Monaghan 形式，电导取 $\rho\kappa$ 的调和平均，写成守恒的质量形式 $m_i\dot c_i=\sum_j2V_iV_j\bar k_{ij}F_{ij}(c_j-c_i)$），与网络用同一组账本密度、同一个隐式时间步，只差空间离散。
\begin{table}[!htbp]\centering\small\caption{末时刻扩散标量的相对 $L^2$ 误差，以及计算时间。}\resizebox{\linewidth}{!}{\begin{tabular}{lccccc}\toprule 算例 & 网络·质量 & SPH & 网络·体积 & 纯对流 & 时间（网络 / SPH）\\\midrule
双环流（$\kappa=10^{-3}$） & 3.2\% & 17.0\% & 2.8\% & 66.1\% & 63 s / 101 s\\
Sod（$\kappa=10^{-3}$） & 0.17\% & 0.88\% & 1.8\% & 7.4\% & 26 s / 48 s\\
爆炸波（$\kappa=2\times10^{-4}$） & 2.2\% & 4.8\% & 22.3\% & 109.1\% & 65 s / 113 s\\
激波–气泡（DNS 标量） & 6.3\% & 6.7\% & 23.6\% & 51.8\% & 15 s / 28 s\\
\bottomrule\end{tabular}}\end{table}
网络质量形式在 Sod、爆炸波、双环流中比 SPH 小 2–5 倍，在激波–气泡中两者相当（参照本身是带数值扩散的 DNS 标量），计算时间约为 SPH 的一半；两者的标量质量都守恒到机器精度，都没有越界。不可压的双环流中体积形式与质量形式等价。SPH 在双环流中误差大，是因为强折叠的物质丝使示踪点分布高度各向异性，各向同性核的 Laplacian 近似变差；P1 刚度对各向异性的星仍然一致。
\begin{figure}[!htbp]\centering\includegraphics[width=\linewidth]{figR_base.pdf}\caption{(a) 密度中位误差：账本、SPH、DTFE；(b) 扩散误差：网络质量形式、SPH、网络体积形式。}\end{figure}
\section{参数敏感性}
默认参数：C 通道阈值 $\tau=\max(0.2,4\hat\sigma)$，锚定周期 5 帧，最多推迟 4 个周期；链式张量 $k=12$。每次只改一个参数。
\begin{table}[!htbp]\centering\small\caption{账本密度的中位误差 / 95 分位，以及 C 通道的检出率 / 误报率。}\resizebox{\linewidth}{!}{\begin{tabular}{lcccc}\toprule 参数 & Sod 密度 & 爆炸波密度 & 激波–涡密度 & 检出 / 误报（Sod；爆炸波）\\\midrule
默认 & $<10^{-6}$ / 0.97\% & 0.58\% / 32.2\% & 0.34\% / 3.4\% & 1.0000 / 1.0e-04；0.9989 / 0.019\\
$\tau_{\min}=0.1$ & $<10^{-6}$ / 0.97\% & 0.58\% / 32.2\% & 0.34\% / 3.5\% & 1.0000 / 2.0e-04；1.0000 / 0.033\\
$\tau_{\min}=0.3$ & $<10^{-6}$ / 0.97\% & 0.59\% / 32.9\% & 0.35\% / 3.4\% & 0.9996 / 6.2e-05；0.9520 / 0.009\\
锚定周期 3 & $<10^{-6}$ / 0.98\% & 0.77\% / 46.3\% & 0.34\% / 2.9\% & 1.0000 / 1.0e-04；0.9989 / 0.019\\
锚定周期 8 & $<10^{-6}$ / 0.91\% & 0.58\% / 32.2\% & 0.35\% / 4.7\% & 1.0000 / 1.0e-04；0.9989 / 0.019\\
不推迟重锚 & $<10^{-6}$ / 1.5\% & 10.5\% / 62.7\% & 0.37\% / 3.5\% & 1.0000 / 1.0e-04；0.9989 / 0.019\\
最多推迟 2 & $<10^{-6}$ / 0.97\% & 2.8\% / 64.0\% & 0.34\% / 3.4\% & 1.0000 / 1.0e-04；0.9989 / 0.019\\
最多推迟 8 & $<10^{-6}$ / 0.97\% & 0.58\% / 32.2\% & 0.34\% / 3.4\% & 1.0000 / 1.0e-04；0.9989 / 0.019\\
\bottomrule\end{tabular}}\end{table}
阈值在所测范围内几乎不影响结果。最敏感的是推迟重锚：不推迟时爆炸波的中位误差从 0.58\% 升到 10.5\%，只允许推迟 2 个周期时为 2.8\%；推迟 4 与 8 个周期结果相同。激波模板持续的帧数超过推迟上限时，锚定星的一端落在激波内，望远相消失效，这与理论一致。锚定周期有取舍：爆炸波中取 3 帧时 95 分位升到 46.3\%，涡核中却偏好短周期（激波–涡的 95 分位在 3、5、8 帧时为 2.9\%、3.4\%、4.7\%）。所以最多推迟次数应按“激波扫过一个星所需帧数 / 锚定周期”设定，锚定周期应按星的畸变程度自适应，而不是固定值。
\begin{table}[!htbp]\centering\small\caption{链式张量估计在爆炸波圆环上的界面平均拉伸（精确值 0.5776）与逐点误差中位数。}\begin{tabular}{lcc}\toprule 参数 & 界面平均 & 逐点误差中位数\\\midrule
默认 & 0.5765 & 0.0091\\
k=8 & 0.5767 & 0.0092\\
k=16 & 0.5763 & 0.0093\\
锚定周期 3 & 0.5766 & 0.0093\\
锚定周期 8 & 0.5765 & 0.0091\\
最多推迟 2 & 0.5767 & 0.0088\\
最多推迟 8 & 0.5765 & 0.0091\\
\bottomrule\end{tabular}\end{table}链式张量对 $k$、锚定周期与推迟次数都不敏感：界面平均在 $\pm2\times10^{-4}$ 内不变。
\section{接近实际数据的检验}
粒子追踪测速（PTV）的位置不确定度通常为 0.05–0.1 像素，示踪点间距约 5–10 像素，即 $\sigma\approx0.01h$–$0.02h$；断轨率取每帧 1\%–3\%。检验包括四项：(i) 每帧独立高斯噪声 $0.02h$、$0.05h$、$0.1h$；(ii) 随机保留 1/4 与 1/16 的示踪点；(iii) 断轨：每条轨迹每帧以概率 $p$ 中断，消失 2 帧后以新身份出现；账本按 5 帧窗口运行，只用窗口内完整的轨迹，新身份的密度由 8 个已知密度的近邻按反距离插值（对数密度）赋初值；(iv) 组合：1/4 稀疏、0.02 倍（稀疏后）间距的噪声、$p=0.01$。
\begin{table}[!htbp]\centering\small\caption{Sod：密度中位误差 / 95 分位。}\resizebox{\linewidth}{!}{\begin{tabular}{lcccccc}\toprule 无扰动 & 噪声 0.02$h$ & 噪声 0.05$h$ & 噪声 0.1$h$ & 稀疏 1/4 & 稀疏 1/16 & 组合\\\midrule
$<10^{-6}$ / 0.97\% & 1.5\% / 5.3\% & 4.1\% / 14.9\% & 9.1\% / 32.4\% & $<10^{-6}$ / 5.6\% & $<10^{-6}$ / 11.7\% & 2.3\% / 14.8\%\\\bottomrule\end{tabular}}
\\[0.3em]\begin{tabular}{lccc}\toprule 断轨概率 $p$ & 0 & 0.01 & 0.03\\\midrule
中位 / 95 分位 & $<10^{-6}$ / 3.0\% & $<10^{-6}$ / 4.9\% & $<10^{-6}$ / 6.9\%\\
末窗口可用轨迹比例 & 100\% & 93\% & 81\%\\\bottomrule\end{tabular}\end{table}
\begin{table}[!htbp]\centering\small\caption{激波–涡（DNS）：密度中位误差 / 95 分位。}\resizebox{\linewidth}{!}{\begin{tabular}{lcccccc}\toprule 无扰动 & 噪声 0.02$h$ & 噪声 0.05$h$ & 噪声 0.1$h$ & 稀疏 1/4 & 稀疏 1/16 & 组合\\\midrule
0.34\% / 3.4\% & 1.9\% / 7.4\% & 4.8\% / 17.3\% & 10.2\% / 35.1\% & 0.85\% / 12.6\% & 1.8\% / 23.7\% & 2.9\% / 11.7\%\\\bottomrule\end{tabular}}
\\[0.3em]\begin{tabular}{lccc}\toprule 断轨概率 $p$ & 0 & 0.01 & 0.03\\\midrule
中位 / 95 分位 & 0.39\% / 3.4\% & 0.41\% / 3.7\% & 0.43\% / 4.1\%\\
末窗口可用轨迹比例 & 100\% & 93\% & 81\%\\\bottomrule\end{tabular}\end{table}
激波–涡涡核区域的位置形式环量与 DNS 的相对误差（末步）：无扰动 3.5\%，噪声 0.02$h$ / 0.05$h$ / 0.1$h$ 为 3.6\% / 3.8\% / 3.9\%，稀疏 1/4 / 1/16 为 11.3\% / 87.4\%。

\paragraph{时间平滑。}对带噪声的位置做 Savitzky–Golay 平滑（二次，窗口 5 或 9 帧）后再算账本，密度中位误差：
\begin{center}\small\begin{tabular}{lcccc}\toprule & \multicolumn{2}{c}{噪声 0.02$h$} & \multicolumn{2}{c}{噪声 0.05$h$}\\ & 不平滑 & 窗 5 / 窗 9 & 不平滑 & 窗 5 / 窗 9\\\midrule
Sod & 1.5\% & 1.4\% / 1.4\% & 4.1\% & 3.5\% / 3.2\%\\
激波–涡 & 2.0\% & 1.8\% / 2.3\% & 4.8\% & 4.1\% / 4.0\%\\
\bottomrule\end{tabular}\end{center}
无噪声数据平滑后（窗 5 / 9）的中位误差为：Sod $<10^{-6}$ / 0.16\%，激波–涡 0.45\% / 1.4\%。平滑只带来 10\%–30\% 的改善，窗口过长还会在激波处抹掉位置轨迹的拐折。
\paragraph{读数。}(1) \textbf{噪声。}账本对每帧独立噪声的响应约与 $\sigma/h$ 成正比：$0.02h$ 时中位误差 1.5\%（Sod）、1.9\%（激波–涡），$0.05h$ 时 4.1\%、4.8\%。按 PTV 的典型精度（$\sigma\approx0.01h$–$0.02h$），中位误差在 1\%–2\% 量级。错帧剖分与有向面积的抗噪版本没有可见改善（$0.05h$ 时 4.8\% 对 4.8\%）。(2) \textbf{稀疏。}示踪点减为 1/4 与 1/16 时，激波–涡的中位误差为 0.85\% 与 1.8\%，涡核区域环量误差为 11.3\% 与 87.4\%。(3) \textbf{断轨。}每帧 1\%–3\% 的断轨下，窗口化账本仍有 81\%–93\% 的轨迹可用，中位误差 0.41\%–0.43\%（激波–涡）；主要代价来自窗口端点落在激波内时损失的望远相消，以及新身份初值的插值误差。(4) \textbf{组合。}三种扰动同时存在时，中位误差 2.3\%（Sod）与 2.9\%（激波–涡），95 分位 14.8\% 与 11.7\%。

\section{复杂应用：单模 Richtmyer–Meshkov}\label{pB:sec:rm}
\paragraph{设置。}二维 Euler（MUSCL–HLLC），$M_s=1.5$ 激波从轻气体（$\rho=1$）射向扰动界面 $x=0.5+a_0\cos(2\pi y)$，$a_0=0.04$，重气体 $\rho=3$（Atwood 数 0.5），半波长通道 $y\in[0,0.5]$（$y=0$ 对称、$y=0.5$ 壁面），$\Delta x=0.005$，$t\in[0,2]$，51 帧；8000 个示踪点（$h=0.004$）在求解器内推进。所有网络量只用示踪点位置与初始密度。
\begin{table}[!htbp]\centering\small\caption{RM：网络结果与 DNS。}
\begin{tabularx}{\linewidth}{@{}L >{\raggedright\arraybackslash}p{0.46\linewidth}@{}}\toprule 量 & 结果\\\midrule
T1 检出率 / 误报（以 DNS 单步密度跃变 15\% 为准） & 1.00 / 0\\
T2 密度中位误差（$t=0.4,1.0,2.0$） & 0.50\%，0.56\%，0.73\%\\
\quad 95 分位（离界面 $>0.03$） & 6.4\%，11.8\%，17.7\%\\
T2 末时刻：账本 / SPH / DTFE 中位误差 & 0.73\% / 4.6\% / 17.1\%\\
T4 界面带环量：位置形式与 DNS 的中位相对误差 & 0.16\%\\
\quad 激波后环量（位置形式 / DNS） & 0.0335 / 0.0335\\
T6 落在平均界面之后的重气体质量：网络 / DNS & $t=0.4$：0.0238 / 0.0206；$t=0.8$：0.0368 / 0.0355；$t=1.2$：0.0526 / 0.0531；$t=1.6$：0.0706 / 0.0725；$t=2$：0.0868 / 0.0890\\
振幅（界面示踪点 / DNS 等值线），5 个时刻 & 相差 $\le$ 1.2\%\\
线性段增长率：实测 / Richtmyer / 由环量 $k\Gamma/4$ & 0.0421 / 0.0433 / 0.0501\\
T9 界面平均：压缩 $-\ln J$ / 拉伸 $\ln\lambda$（$t=2$） & 0.798 / 0.054\\
T9 $\ln\langle\lambda\rangle$ 与 DNS 界面长度 $\ln(L/L_0)$ 的增长（$t=0.4\to2$） & 0.0768 / 0.0757\\\bottomrule
\end{tabularx}\end{table}
\begin{figure}[!htbp]\centering\includegraphics[width=\linewidth]{figR_rm.pdf}
\caption{单模 RM。(a1–a3) DNS 密度；(a4) 末时刻账本密度误差（对数）；(b) 振幅：界面示踪点、DNS 等值线、Richtmyer 模型（参数取自位置）与由沉积环量估计的斜率；(c) 界面物质带环量，位置形式与 DNS；(d) 界面掺混归因。}\end{figure}

\paragraph{物理读数。}
(1) \textbf{增长率。}激波扫过（$t\approx0.16$）后，界面示踪点给出的振幅先被压缩到 $a_0^+=0.0275$，随后线性增长，$t\in[0.6,1.4]$ 的增长率为 0.0421，$t>1.4$ 降到 0.0389（开始非线性饱和）。Richtmyer 冲击模型 $\dot a=k\,\Delta u\,A^+a_0^+$ 给出 0.0433，与实测相差 3\%；其中界面速度跳跃 $\Delta u=0.510$、波后 Atwood 数 $A^+=0.491$（由账本密度 2.17 与 6.37 求得）、$a_0^+$ 全部取自示踪点位置。
(2) \textbf{环量与增长。}把沉积在半波长界面带上的环量看成正弦涡片，$\dot a=k\Gamma/4$ 给出 0.0501，比实测高约 19\%：这一估计假设涡片无限薄，并忽略界面带内非涡片部分的涡量。位置形式的环量与 DNS 中位误差为 0.16\%，所以这一偏差来自模型而不是测量。
(3) \textbf{掺混归因。}到 $t=2$（$ka\approx0.6$）界面压缩为 0.80，几乎全部在激波扫过时一次获得；拉伸从接近零缓慢增长到 0.054。所以在这一阶段，单模 RM 界面上的梯度增长约 94\% 来自压缩；拉伸只有在卷起（$ka>1$）之后才可能接管。此时尚未卷起，界面带也没有出现激波–重气泡 $t=1.4$ 时的不收敛问题。
\section{统一数值登记表}
四份旧报告（质量输运报告、界面拉伸重推报告、六算例统一报告、T1–T9 定稿）中同一量有不同数值，或已被推翻的结论。下表给出以后引用时唯一采用的版本。
\begin{table}[!htbp]\centering\small\caption{数值登记。“现行值”为本稿或最近一次统一流水线的结果。}
\begin{tabularx}{\linewidth}{@{}>{\raggedright\arraybackslash}p{0.24\linewidth} >{\raggedright\arraybackslash}p{0.27\linewidth} L@{}}\toprule 量 & 旧值（出处） & 现行值与说明\\\midrule
激波–气泡界面拉伸，$t=1.4$ & 0.60–1.05 区间（质量输运报告） & 0.285（链式张量），参照 0.390，涡片范围 [0.06, 0.39]；旧区间来自点对与折线法，已推翻\\
激波–气泡压缩占总增长，$t=1.4$ & 30\%–43\%（质量输运报告）；54\%–62\%（重推报告） & 依赖界面带的压缩值，DNS 加密后不收敛，降级为待粘性 DNS 复核；$t\le0.6$ 的“先切向缩短、抵消大部分压缩”稳健\\
激波–气泡界面长度比 & 5.3（质量输运报告） & 2.16（自适应标记曲线）\\
激波–气泡界面压缩，$t=1.4$ & 0.450（重推报告）；0.501（统一报告） & 不确定：三档 DNS 轨迹上账本为 0.457 / 0.468 / -0.109，DNS 为 0.416 / 0.363 / 0.358；卷起后的界面带处于信息极限，见第~\ref{pB:sec:dns} 节\\
Lamb–Oseen 圆环拉伸 & “约为 0”（质量输运报告） & $-0.31$（逐点精确值的界面平均）\\
激波–涡梯度能量占比 & 0.75–0.81（统一报告） & 0.53（Kabsch 对齐后）；原始值约三分之一来自平移，不作物理读数\\
Sod 冲击部分 / 总变化 & 471\%（统一报告） & 删除（分母接近零）\\
激波–气泡 Dirac 残差（激波处） & 0.005（正文）/ 0.01（总表）（统一报告） & 0.005；T5 不进入论文 A\\
密度中位误差：双环流 / 激波–涡 / 激波–气泡 & 1.3\% / 0.33\% / 0.70\%（质量输运）；1.35\% / 0.36\% / 0.73\%（统一） & 1.3\% / 0.34\% / 0.67\%（本稿，只凭位置的热步，DNS 网格 0.005）\\
激波–气泡扩散误差（质量 / 体积） & 6.1\% / 24\%（质量输运）；5.1\% / 25.2\%（统一） & 6.3\% / 23.6\%（本稿；SPH 6.7\%）\\
T8 方差衰减的条件 & “需 $K_{ij}\le0$，Delaunay 满足” & 方差衰减无条件；$K_{ij}\le0$ 只关系离散极值原理，变电导下 Delaunay 不保证\\
位置形式 Kelvin 的边向量 & $\tfrac14(q^{n-1}+2q^n+q^{n+1})$ & 一族精确形式，$c=0$ 即 $q^n$ 最简\\\bottomrule
\end{tabularx}\end{table}
\section{仍然存在的缺口}
\begin{enumerate}
\item \textbf{卷起后的界面。}激波–重气泡 $t=1.4$ 界面带上的压缩，以及依赖它的掺混占比，需要在有物理粘性、界面网格收敛的 DNS 上复核，且示踪点间距要不大于最薄的物质层。
\item \textbf{自适应锚定。}锚定周期在爆炸波中宜长、在涡核中宜短；应按锚定星的畸变程度（如累积 $|\mu|$）决定重锚，而不是固定帧数。
\item \textbf{三维。}没有三维算例。理论推广直接（四面体、三维 Nanson），可作为局限写明。
\item \textbf{真实实验数据。}本稿的“实际数据”检验是在 DNS 轨迹上加扰动，仍不是真实 PTV 数据；SBLI 需要粘性求解器与大规模示踪点，本稿以单模 RM 代替。
\item \textbf{噪声。}账本对每帧独立噪声的响应约为 $\sigma/h$ 量级；抗噪版本（错帧剖分、有向面积）没有明显改善，时间平滑只改善 10\%–30\%。
\item \textbf{断轨。}窗口化账本在断轨下可用，但窗口端点落在激波内时会损失望远相消；需要让窗口端点也按激波模板推迟。
\item \textbf{文献定位}仍需补查：free-Lagrange 方法、Froyland \& Junge（2018）有限元动态 Laplacian、DEC 与图形学中的 Helmholtz–Hodge 分解。
\end{enumerate}

\clearpage\part{第三部分　缺口补齐：粘性 DNS、自适应锚定、噪声传播律、动态 Laplacian、环量沉积、文献定位}

\noindent\fcolorbox{accent!40}{tint}{\begin{minipage}{0.965\linewidth}\vspace{0.3em}
\textbf{\sffamily 结论}\quad (1) \textbf{粘性 DNS}：界面带收敛：两档细网格的账本相差 0.009、DNS 相差 0.004，最细网格上二者相差 0.016；$t=1.4$ 拉伸 -0.019、压缩 0.322（$Re\approx1300$）。(2) \textbf{自适应锚定}：激波–涡 95 分位 3.4\% → 2.9\%，其余算例不变差，解决了锚定周期的取舍。(3) \textbf{噪声传播律}：账本中位误差 $\approx0.74\,\sigma/h$，只由两端帧决定、不随累积变形放大，静态与流动数据都吻合。(4) \textbf{T7 改为加权动态 Laplacian}：Lamb–Oseen 涡核 Jaccard 0.11 → 0.82；双环流 0.59 → 0.67；它与 T8 共用同一算子。(4$'$) \textbf{环量沉积}：只凭位置的上半气泡环量 0.266，与 DNS（0.270）相差 1.5\%，与 Samtaney–Zabusky 的 SZ1 模型（0.289）相差 8\%。(5) \textbf{文献}：附件中 Froyland–Junge（2018）与 Froyland–Kwok（2020）已经给出了轨迹上的 FEM 刚度与加权动态 Laplacian，Karrasch–Keller（2020）给出了拉格朗日热方程；论文 A 的新意应改述为“只凭位置得到可压缩、含激波流动的质量测度（账本），并据此求解变密度的拉格朗日对流—扩散与掺混归因”，而不是 FEM 离散本身。
\vspace{0.3em}\end{minipage}}



\section{缺口与处理一览}
\begin{table}[!htbp]\centering\small\caption{上一份报告列出的缺口与本轮的处理。}
\begin{tabularx}{\linewidth}{@{}>{\raggedright\arraybackslash}p{0.2\linewidth} >{\raggedright\arraybackslash}p{0.3\linewidth} L@{}}\toprule
缺口 & 本轮的处理 & 结果\\\midrule
卷起后界面带的压缩参照 & 给求解器加物理粘性、热传导与组分扩散（$\mu=5\times10^{-4}$，$Re\approx1300$），三档网格 & 界面带收敛：两档细网格的账本相差 0.009、DNS 相差 0.004，最细网格上二者相差 0.016；$t=1.4$ 拉伸 -0.019、压缩 0.322（$Re\approx1300$）\\
锚定周期的取舍 & 残差触发的自适应锚定：星的仿射拟合残差超过阈值即提前重锚 & 激波–涡 95 分位 3.4\% → 2.9\%，其余算例不变差\\
噪声没有理论 & 推导账本的噪声传播律 & 中位误差 $\approx0.74\,\sigma/h$，与静态与流动数据一致\\
T7 结果弱 & 用加权 FEM 动态 Laplacian（Froyland–Junge 2018 的离散，Froyland–Kwok 2020 的权重）取代 Ulam 矩阵，质量取自账本 & Lamb–Oseen 涡核 Jaccard 0.11 → 0.82；双环流 0.59 → 0.67\\
环量沉积的模型比较 & 按 Samtaney–Zabusky（1994）式 4.12 与 5.12 手工录入公式，与位置形式环量比较 & 上半气泡环量：位置形式 0.266、DNS 0.270、SZ1 模型 0.289（相差约 8\%）\\
文献定位 & 阅读附件 15 篇，逐篇对应 & 见第~\ref{pC:sec:lit} 节；T8 的 FEM 刚度与加权动态 Laplacian 已有先例，论文 A 的新意需改述\\
断轨窗口端点 & 未做 & 方案见第~\ref{pC:sec:rest} 节\\
三维、真实 PTV & 未做 & 作为局限写明\\\bottomrule
\end{tabularx}\end{table}
\section{粘性 DNS：卷起之后的界面压缩}\label{pC:sec:visc}
\paragraph{设置。}在原 MUSCL–HLLC 求解器上加守恒形式的粘性、热传导与组分扩散通量（常数 $\mu=5\times10^{-4}$，$Pr=0.72$，$\rho D=\mu$），其余设置与原激波–重气泡相同（$M_s=1.5$、密度比 3），示踪点不变（$h=0.005$）。以气泡直径与波后速度计，$Re\approx1300$。界面的扩散厚度 $\sqrt{\nu t}$ 在 $t=1.4$ 约为 0.015，即 3 个示踪点间距，因此物质层不再薄于 $h$。网格取 $\Delta x=0.01,0.005,0.0035$。
\begin{table}[!htbp]\centering\small\caption{粘性激波–重气泡，$t=0.6$：各带的账本压缩 / DNS 插值压缩。}
\begin{tabular}{lccc}\toprule
$\phi_0$ 区间 & $\Delta x=0.01$ & $\Delta x=0.005$ & $\Delta x=0.0035$\\\midrule
$[-0.06,-0.03]$ & 0.568 / 0.564 & 0.580 / 0.590 & 0.583 / 0.593\\
$[-0.03,-0.015]$ & 0.420 / 0.368 & 0.425 / 0.402 & 0.424 / 0.414\\
$[-0.015,-0.005]$ & 0.291 / 0.222 & 0.210 / 0.158 & 0.186 / 0.153\\
$[-0.003,0.003]$ & 0.363 / 0.321 & 0.320 / 0.286 & 0.316 / 0.281\\
$[0.005,0.015]$ & 0.605 / 0.655 & 0.723 / 0.781 & 0.756 / 0.797\\
$[0.015,0.03]$ & 0.651 / 0.699 & 0.666 / 0.685 & 0.669 / 0.678\\
$[0.03,0.06]$ & 0.557 / 0.583 & 0.553 / 0.562 & 0.553 / 0.557\\
\bottomrule\end{tabular}\end{table}
\begin{table}[!htbp]\centering\small\caption{粘性激波–重气泡，$t=1.4$：各带的账本压缩 / DNS 插值压缩。}
\begin{tabular}{lccc}\toprule
$\phi_0$ 区间 & $\Delta x=0.01$ & $\Delta x=0.005$ & $\Delta x=0.0035$\\\midrule
$[-0.06,-0.03]$ & 0.485 / 0.433 & 0.486 / 0.481 & 0.490 / 0.493\\
$[-0.03,-0.015]$ & 0.315 / 0.234 & 0.279 / 0.243 & 0.266 / 0.247\\
$[-0.015,-0.005]$ & 0.195 / 0.146 & 0.092 / 0.050 & 0.066 / 0.036\\
$[-0.003,0.003]$ & 0.333 / 0.352 & 0.312 / 0.310 & 0.322 / 0.306\\
$[0.005,0.015]$ & 0.647 / 0.752 & 0.794 / 0.868 & 0.834 / 0.881\\
$[0.015,0.03]$ & 0.750 / 0.834 & 0.791 / 0.814 & 0.798 / 0.808\\
$[0.03,0.06]$ & 0.676 / 0.709 & 0.681 / 0.690 & 0.682 / 0.687\\
\bottomrule\end{tabular}\end{table}
\begin{table}[!htbp]\centering\small\caption{粘性激波–重气泡：密度误差与 T9 掺混归因随网格的变化。}
\begin{tabular}{lccc}\toprule 量 & $\Delta x=0.01$ & $\Delta x=0.005$ & $\Delta x=0.0035$\\\midrule
$t=0.6$ 密度中位 / 95 分位 & 0.51\% / 8.9\% & 0.29\% / 5.9\% & 0.27\% / 4.5\%\\
$t=1.4$ 密度中位 / 95 分位 & 1.1\% / 15.6\% & 0.53\% / 8.9\% & 0.48\% / 7.1\%\\
$t=0.6$ 拉伸 / 压缩 / 总增长 & -0.331 / 0.363 / 0.031 & -0.331 / 0.320 / -0.011 & -0.330 / 0.316 / -0.014\\
$t=1$ 拉伸 / 压缩 / 总增长 & -0.349 / 0.402 / 0.053 & -0.346 / 0.374 / 0.028 & -0.335 / 0.378 / 0.044\\
$t=1.4$ 拉伸 / 压缩 / 总增长 & -0.036 / 0.333 / 0.297 & -0.029 / 0.312 / 0.283 & -0.019 / 0.322 / 0.302\\
$x>1$ 的重气体质量：网络 / DNS & 0.2864 / 0.2889 & 0.2888 / 0.2916 & 0.2892 / 0.2925\\\bottomrule\end{tabular}\end{table}
\paragraph{读数。}在 $|\phi_0|<0.003$ 的界面带上，$t=1.4$ 时账本压缩随网格为 0.333 / 0.312 / 0.322，DNS 为 0.352 / 0.310 / 0.306；两档细网格之间账本相差 0.009、DNS 相差 0.004，最细网格上二者相差 0.016。有了物理粘性，界面带上账本与 DNS 都随网格收敛并且彼此接近，所以无粘 DNS 中的不收敛来自没有物理耗散的 Kelvin–Helmholtz 小尺度结构（细于示踪点间距），而不是账本本身。 掺混归因：$t=1.4$ 时界面平均拉伸随网格为 -0.036 / -0.029 / -0.019，压缩为 0.333 / 0.312 / 0.322，梯度总增长为 0.297 / 0.283 / 0.302。拉伸接近零而略为负，所以在 $Re\approx1300$ 下，到 $t=1.4$ 界面梯度的增长几乎全部来自压缩：物理粘性抑制了卷起，界面没有被明显拉长。这是另一个雷诺数下的结论，不能直接替代无粘算例中被降级的“54\%–62\%”；它说明压缩与拉伸的分配依赖雷诺数，而账本在两个雷诺数下都给出了可检验的值（拉伸与压缩接近抵消时，不宜再用“压缩占总增长的比例”作指标，应分别报告两项）。
\section{自适应锚定}
\paragraph{规则。}账本精确的条件是：锚定星的顶点在锚定段内作仿射运动（星多边形的边保持为物质直线）。于是每一步对每个示踪点的锚定星做最小二乘仿射拟合，残差 $r_i=\big(\tfrac1K\sum_k|x_k-Fx^a_k-c|^2\big)^{1/2}/\sqrt{A^\ast_i}$；当 $r_i>\varepsilon$ 或锚定段达到 5 帧时重锚，但星触及激波模板时照旧推迟（总长不超过 20 帧）。这样，仿射性好的区域保持 5 帧的锚定，涡核等非仿射区域自动缩短锚定段。
\begin{table}[!htbp]\centering\small\caption{密度中位误差 / 95 分位：固定 5 帧与自适应锚定（$\varepsilon=0.01,0.02,0.04$）。}\resizebox{\linewidth}{!}{\begin{tabular}{lcccc}\toprule 算例 & 固定 5 帧 & $\varepsilon=0.01$ & $\varepsilon=0.02$ & $\varepsilon=0.04$\\\midrule
Sod & $<10^{-6}$ / 0.97\% & $<10^{-6}$ / 0.96\% & $<10^{-6}$ / 0.97\% & $<10^{-6}$ / 0.97\%\\
爆炸波 & 0.58\% / 32.2\% & 0.58\% / 32.2\% & 0.58\% / 32.2\% & 0.58\% / 32.2\%\\
双环流 & 1.3\% / 5.7\% & 1.2\% / 5.0\% & 1.2\% / 5.4\% & 1.2\% / 5.7\%\\
Lamb–Oseen & 0.41\% / 13.0\% & 0.40\% / 7.1\% & 0.40\% / 8.8\% & 0.41\% / 10.7\%\\
激波–涡 & 0.34\% / 3.4\% & 0.34\% / 2.9\% & 0.34\% / 3.3\% & 0.35\% / 3.4\%\\
激波–气泡 & 0.67\% / 29.7\% & 0.68\% / 29.5\% & 0.67\% / 29.3\% & 0.67\% / 29.5\%\\
RM & 0.73\% / 35.4\% & 0.73\% / 35.3\% & 0.73\% / 35.4\% & 0.73\% / 35.3\%\\
\bottomrule\end{tabular}}\end{table}
自适应锚定在所有算例上都不比固定 5 帧差，在涡核主导的激波–涡中把 95 分位从 3.4\% 降到 2.9\%（$\varepsilon=0.01$），与“锚定周期 3 帧”的最好结果相当，同时保留了爆炸波中长锚定的好处（32.2\%，与固定 5 帧相同）。上一轮的“锚定周期需取舍”由此解决；建议默认 $\varepsilon=0.01$。早先试过的“只靠残差触发、最长 20 帧”的版本在爆炸波中变差（95 分位 39.6\%），原因是光滑膨胀区的锚定段过长，这说明 5 帧上限本身有用：它限制了单段内非仿射误差的积累。
\section{噪声传播律}
\begin{proposition}[账本的噪声传播]\status{已证+数值}
设每帧位置带独立各向同性噪声 $\delta x\sim\mathcal N(0,\sigma^2I)$。(i) 节点的星面积就是其邻点组成的连接多边形的面积，与中心点位置无关；对邻点 $x_k$（按角度排序）有 $\partial A^\ast/\partial x_k=\tfrac12R_{90^\circ}(x_{k+1}-x_{k-1})$，故
\begin{equation}
\operatorname{std}(\delta\ln A^\ast)=\frac{\sigma}{2A^\ast}\Big(\sum_k|x_{k+1}-x_{k-1}|^2\Big)^{1/2}=c_\ast\,\frac{\sigma}{h},
\end{equation}
正六边形星 $c_\ast=\sqrt{4.5}/(3\sqrt3/2)=0.816$。(ii) 账本 $\ln J$ 的误差主要来自窗口两端的两帧：中间重锚时，旧星在当前帧的面积与新星在同一帧的锚定面积用的是同一组带噪位置，且两星大部分顶点相同，噪声大部分相消。于是 $\operatorname{std}(\delta\ln J)\approx\sqrt2\,c_\ast\sigma/h$，中位绝对误差 $\approx0.674\sqrt2\,c_\ast\,\sigma/h$。
\end{proposition}
\paragraph{检验。}在静止的抖动格点上（无流动，只有噪声），由实际剖分算得 $c_\ast=0.777$，预测中位误差 $0.741\,\sigma/h$；账本实测为：11 帧 $0.73$–$0.79\,\sigma/h$，41 帧 $0.80$–$0.89\,\sigma/h$（随帧数只缓慢增长，印证了 (ii) 的相消）。流动算例中，Sod 为 $0.75\,\sigma/h$，激波–涡为 $1.0\,\sigma/h$（涡核的非仿射误差叠加在噪声上）。

\paragraph{推论。}(1) 要使密度中位误差低于 1\%，需要 $\sigma\lesssim0.013h$，与 PTV 的典型精度（$0.01h$–$0.02h$）同量级；(2) 误差按局部间距计：压缩区间距变小，同样的 $\sigma$ 带来更大的相对误差，压缩比 $\eta$ 时放大约 $\sqrt\eta$；(3) 与基于速度梯度的诊断不同（Badza 等 2023 与 Balasuriya 2020 讨论的不确定性随拉伸放大），账本的噪声误差不随累积变形放大；随变形放大的是链式张量估计（放大因子为各向异性 $A$）。(4) 噪声主要在两端帧，所以对首末各若干帧的位置做时间平均，比全程平滑更有效。
\section{用加权动态 Laplacian 取代 T7}
T7 的按质量加权 Ulam 矩阵在六个算例中三个退化、一个找错主导集（Lamb–Oseen 与涡核的 Jaccard 只有 0.11）。附件中的 Froyland \& Junge（2018）给出了散乱轨迹上的 FEM 动态 Laplacian：每个时刻对当前位置做 Delaunay 剖分，组装 P1 刚度并对时间平均，解广义特征问题；Froyland \& Kwok（2020）把它推广到非保体积的动力学与非均匀初始测度。这两者与论文 A 的算子完全对应：P1 刚度就是 T8 的 $K(t)$，权重（测度）就是账本给出的质量。所以 T7 应改为
\begin{equation}
\bar K v=\lambda M v,\qquad \bar K=\frac1S\sum_s K(t_s),\quad M=\mathrm{diag}(m_i),
\end{equation}
可压缩时 $K(t_s)$ 以账本密度加权。
\paragraph{检验。}Lamb–Oseen 涡对：第 2–4 个特征向量取 90\% 分位以上的集合，与涡核（$r_0<1.5a$）的最佳 Jaccard 为 0.82（Ulam 为 0.11）。双环流：第二特征向量按中位数二分，与左环流的 Jaccard 为 0.67（Ulam 为 0.59）。所以 T7 改为动态 Laplacian 后不再是弱项，而且它与 T8 共用同一个算子，与第 1 节的“质量度量下的 Dirac 算子”框架一致：$\bar K$ 是 $D^2$ 节点块的时间平均。需要明确的是，这一部分的离散方法来自 Froyland–Junge，权重的理论来自 Froyland–Kwok；论文 A 的新内容是用账本从位置得到可压缩流中的权重。
\section{环量沉积：与 Samtaney--Zabusky 模型比较}\label{pC:sec:sz}
Samtaney \& Zabusky（1994，JFM 269）给出平面界面单位原长度上的沉积环量（其式 4.12，换算到以入射气体声速 $c_0$ 为单位）
\begin{equation}
\Gamma'\approx\frac{2}{\gamma+1}\Big(1-\eta^{-1/2}\Big)\sin\alpha\,\big(1+M^{-1}+2M^{-2}\big)(M-1)\,c_0,
\end{equation}
以及圆形气泡（半个气泡）上的总环量模型 SZ1：$\Gamma_{iii}=(1+\tfrac12\pi)\Gamma'_1r_0$（其式 5.12，$\Gamma'_1$ 为上式 $\sin\alpha$ 的系数），第二阶段末约为 $\Gamma'_1r_0$。本算例 $M=1.5$、$\eta=3$、$\gamma=1.4$、$r_0=0.25$、$c_0=1$，得 $\Gamma'_1=0.4500$，第二阶段末 0.1125，SZ1 为 0.2892。

位置形式的环量取初始距界面 $|\phi_0|<0.15$ 的物质带（上半气泡），由中点位移上链的环量求和得到（T4，只凭位置）。
\begin{table}[!htbp]\centering\small\caption{上半气泡的环量（取绝对值）。}\begin{tabular}{lccc}\toprule & $t=0.6$ & $t=1.0$ & $t=1.4$\\\midrule
无粘，$\Delta x=0.005$：位置形式（物质带） & 0.266 & 0.269 & 0.274\\
\quad DNS，同一组三角形上积分 & 0.264 & 0.296 & 0.430\\
\quad DNS，整个上半区域 & 0.270 & 0.265 & 0.242\\
无粘，$\Delta x=0.0025$：位置形式（物质带） & 0.265 & 0.263 & 0.268\\
\quad DNS，同一组三角形上积分 & 0.261 & 0.195 & 0.574\\
\quad DNS，整个上半区域 & 0.270 & 0.264 & 0.242\\
粘性，$\Delta x=0.0035$：位置形式（物质带） & 0.265 & 0.264 & 0.264\\
\quad DNS，同一组三角形上积分 & 0.265 & 0.263 & 0.262\\
\quad DNS，整个上半区域 & 0.270 & 0.261 & 0.240\\
SZ1 模型（第三阶段末） & \multicolumn{3}{c}{0.289}\\\bottomrule\end{tabular}\end{table}
\paragraph{读数。}(1) 激波扫过后（$t=0.6$），只凭位置得到的上半气泡环量为 0.266，DNS 整个上半区域为 0.270，二者相差约 1.5\%；SZ1 模型为 0.289，比两者高约 8\%，在 Samtaney–Zabusky 报告的模型与 DNS 的偏差范围内。所以 T4 的位置形式环量不仅与本文 DNS 一致，也与经典模型一致。(2) 无粘算例中，物质带的环量在 $t=0.6$ 之后几乎守恒（Kelvin 定理，0.266 → 0.274），而“同一组三角形上用顶点涡量积分”的 DNS 值在后期偏离（$t=1.4$ 为 0.430）：卷起后物质三角形严重变形并跨过涡片，顶点平均的求积失效，而位置形式是多边形边界上的精确环量，不受影响。(3) 粘性算例中两者在全程一致（中位相对误差 0.35\%），窄带（$|\phi_0|<0.04$）的环量随时间减小，宽带不变，这正是粘性把涡量扩散出窄带、而宽带的边界上粘性通量很小的结果。
\section{附件文献的定位}\label{pC:sec:lit}
\begin{table}[!htbp]\centering\small\caption{附件文献与论文 A 的关系。}
\begin{tabularx}{\linewidth}{@{}>{\raggedright\arraybackslash}p{0.26\linewidth} L@{}}\toprule
文献 & 与本工作的关系与应如何引用\\\midrule
Froyland \& Junge 2018，SIADS（散乱、稀疏、不完整轨迹上的 FEM 动态 Laplacian） & 逐时刻 Delaunay 剖分、P1 刚度、时间平均，正是 T8 刚度与新 T7 的离散；论文 A 不能再声称“轨迹上的 P1 刚度”是新的。新的是：瞬态对流—扩散的求解、以 $\rho\kappa$ 为电导、以账本质量为质量矩阵，以及在激波流中对 DNS 的定量检验\\
Froyland \& Kwok 2020，JNS（加权黎曼流形上的动态 Laplacian） & 非保体积动力学与非均匀初始测度的理论；“度量即测度”在连续层面已由它给出。论文 A 的贡献是在可压缩、含激波的流动中只凭位置得到这一测度，并定量显示面积权重的误差（5–10 倍）\\
Karrasch \& Keller 2020（LCS 的几何热流理论） & 拉格朗日坐标下的对流—扩散是时变热方程，时间平均即动态 Laplacian；T8 的 $M\dot c=-K(t)c$ 是它的 FEM 离散，并推广到变密度与 $\rho\kappa$\\
Froyland 2015（动态等周） & 动态 Laplacian 的出发点；新 T7 的理论依据\\
Hadjighasem 等 2017（LCS 方法的比较） & 脊分类工作的基准与比较对象；不进入论文 A\\
Haller 等 2021（单轨迹准客观诊断）；Aksamit 等 2024（稀疏轨迹的相对拉伸与转动） & “只凭轨迹”的已有路线：前者用单条轨迹，后者用相对量；AFRC 椭圆通道直接用了后者。论文 A 应说明其区别：三角形上的有限步恒等式给出严格的积分量（面积、环量），而不是逐点的准客观量\\
Farazmand \& Haller 2016（极转角 PRA）；Haller \& Beron-Vera 2013（黑洞涡） & 三点映射的 $\arg p$ 即三角形上的极转角，合并块给出跨激波的 PRA；椭圆 LCS 的参照\\
Badza 等 2023；Balasuriya 2020（随机敏感度） & 速度不确定性下 LCS 的敏感度；第 4 节的噪声传播律与之对照：账本误差 $O(\sigma/h)$，不随累积变形放大\\
Singer \& Wu 2012（向量扩散映射与联络 Laplacian） & 以局部旋转为联络的扩散；可用于把逐步的 Kabsch 对齐推广为局部客观的联络，属于拓扑信号那一篇\\
Samtaney \& Zabusky 1994；Picone \& Boris 1988；Haas \& Sturtevant 1987 & 激波–气泡的环量沉积模型与实验。论文 A 的位置形式环量应与 Samtaney–Zabusky 的平面界面标度律（$\Gamma'\propto(1-\eta^{-1/2})\sin\alpha\,(1+M^{-1}+2M^{-2})(M-1)$）与气泡总环量模型比较；本轮已做，见第~\ref{pC:sec:sz} 节\\
Anand \& Bianconi 2009；Anand, Bianconi \& Severini 2011；多路网络熵与重叠；带味网络几何；《多层网络》 & 网络系综的 Shannon 熵与以 $\rho=L/\operatorname{Tr}L$ 定义的 von Neumann 熵。可用于拓扑信号那一篇：以质量度量下的 $\bar K$ 定义密度矩阵，其 von Neumann 熵度量“混合几何”的谱复杂度；时间窗层可用多路网络熵与重叠刻画相干结构的存续。不进入论文 A\\\bottomrule
\end{tabularx}\end{table}
\section{仍未完成}\label{pC:sec:rest}
\begin{enumerate}
\item \textbf{断轨窗口端点。}方案：不再把账本切成固定窗口，而是让每条轨迹在其连续段内保持自己的锚定星；当锚定星的某个顶点消失时，若不在激波模板内则立即重锚，否则推迟到激波过后、用剩余顶点的连接多边形计算面积比。需要重写账本的数据结构（每帧的点集不同），本轮未做。
\item \textbf{三维与真实 PTV 数据}：作为论文 A 的局限写明。
\end{enumerate}

\clearpage\part{第四部分　断轨下的逐轨账本}

\noindent\fcolorbox{accent!40}{tint}{\begin{minipage}{0.965\linewidth}\vspace{0.3em}
\textbf{\sffamily 结论}\quad 逐轨账本把覆盖率从窗口化的 79\%–95\% 提高到 100\%，并在大多数断轨情形下降低了误差；在“粒子恰在穿越激波时断轨”的对抗情形中，Sod 的 95 分位由窗口化的 5.6\% 降到 1.9\%。关键设计是：顶点消失时缩减环而不中断测量；环在激波中断开时，在锚定帧用幸存轨迹重建环，使锚定点留在激波前。
\vspace{0.3em}\end{minipage}}



\section{问题}
PTV 数据中的轨迹会中断：粒子在若干帧内丢失，之后以新的身份重新出现。上一版的处理是把账本切成固定的 5 帧窗口，只用窗口内完整的轨迹，新身份的密度由近邻插值赋值。它有两个缺陷：(1) 窗口端点若落在激波内，锚定星的一端就在激波中，望远相消失效；(2) 只要轨迹在窗口内断过一次，这一窗口就不能用，覆盖率随断轨率迅速下降。本稿按上一份报告给出的方案，把账本改成逐轨的：每条轨迹在其连续段内保持自己的锚定星。

\section{设计}
\subsection{数据结构}
\begin{itemize}
\item \textbf{轨迹表}：$P\in\mathbb R^{S\times M\times2}$，第 $m$ 列是第 $m$ 条轨迹，缺失处为 NaN。同一个物理粒子可以对应多条轨迹，算法不知道这种对应关系。
\item \textbf{每条轨迹的锚定状态}：锚定帧 $a_i$；锚定时刻 Delaunay 星的连接环（按角度排序的邻居轨迹编号 $r_{i,1},\dots,r_{i,K}$）；这些邻居相对 $i$ 的锚定位置；每对相邻环顶点是否构成三角形的标志；当前段的起点值 $b_i=\ln J_i(a_i)$。
\item \textbf{密度标定}：第 0 帧存在的轨迹用已知的初始密度；之后出生的轨迹单独标定（见下）。
\end{itemize}

\subsection{规则}
\begin{enumerate}
\item \textbf{测量。}$\ln J_i(t)=b_i+\ln\big(A_i^{\rm ring}(t)/A_i^{\rm ring}(a_i)\big)$，其中 $A^{\rm ring}$ 是以 $i$ 为中心、只取当前仍然存在的环顶点的有向扇形面积；分子分母用\emph{同一组}顶点。一个顶点消失时，相邻两个三角形并成一个（用弦代替两条边），而不是中断测量。
\item \textbf{常规重锚}：锚定段达到 5 帧，或环的仿射拟合残差超过 0.01（自适应锚定），重锚到当前帧的星；环触及激波模板时推迟，最多 15 帧。
\item \textbf{激波中的断环。}若存在的环顶点少于 3 个或面积不足锚定时的 30\%，而此时环正处在激波模板内，就\emph{在锚定帧重建环}：取在锚定帧与当前帧都存在的轨迹，在锚定帧做局部 Delaunay，得到新环；锚定帧与 $b_i$ 不变。这样锚定星仍在激波前，穿越照样望远相消。不在激波中的断环直接在当前帧重锚。
\item \textbf{出生与标定}：出生的轨迹从出生帧开始自己的账本（$\ln J=0$）。它的绝对密度由已标定的近邻在标定时刻 $t_c$ 的 $\ln\rho$ 按反距离平均得到，再由自己的账本推到其他时刻。$t_c$ 取出生后第一个自己不在激波模板内的帧（\textbf{延迟标定}），或直接取出生帧（\textbf{即时标定}）。
\end{enumerate}

\begin{proposition}[缩减环的精确性]\status{已证}
若环在锚定段内作仿射运动，则以任意一组存在的环顶点（$\ge3$ 个、按原角度顺序）构成的多边形，其面积比都等于 $\det F$。因此顶点消失不引入误差，只缩小了用于测量的物质多边形。平面激波完整扫过锚定星时，锚定帧与当前帧之间的映射是仿射的，所以在锚定帧重建的环（只用幸存轨迹）同样给出精确的跨激波面积比。
\end{proposition}
\begin{proof}
仿射映射把多边形映为多边形，面积乘以 $\det F$，与顶点个数无关。第二句：锚定帧在激波前、当前帧在激波后时，两帧之间整个邻域经历同一个单轴映射。
\end{proof}

\paragraph{开发中的一个教训。}第一版在“激波中的断环”时也在当前帧重锚。这样锚定被放在了激波中间，在 30\% 粒子于激波处断轨的情形下，Sod 的末帧 95 分位达到 32\%；新出生轨迹的误差也随之变大，因为它们从出错的邻居那里标定。改为“在锚定帧重建环”后，降到 2.0\%（第~\ref{pD:sec:res} 节）。这说明逐轨账本的关键不在缩减环本身，而在不让断轨把锚定点拖进激波。

\section{检验设计}
\begin{itemize}
\item \textbf{算例}：Sod（精确轨迹与密度，7500 个粒子，41 帧）、柱面爆炸波（精确，9345 个，61 帧）、激波–等熵涡（DNS 在线示踪点，14400 个，31 帧）。
\item \textbf{断轨模型}：每帧之后，每条存活轨迹以概率 $p$ 中断；粒子随后缺失 1–3 帧（均匀随机），再以新编号出现。$p=1\%,3\%,5\%$。
\item \textbf{对抗情形}：30\% 的被激波扫过的粒子恰在穿越激波的那一步断轨（外加 1\%/帧的随机断轨），专门检验“端点落在激波内”。
\item \textbf{组合情形}：随机保留 1/4 的粒子，加 0.02 倍（稀疏后）间距的位置噪声，$p=3\%$。
\item \textbf{比较}：窗口化账本（旧）、逐轨账本加即时标定、逐轨账本加延迟标定；以及同一组粒子不断轨时的账本（下限）。
\item \textbf{指标}：末帧仍存在的轨迹中能给出密度的比例（覆盖率）；末帧密度相对误差的中位数与 95 分位；其中出生轨迹的 95 分位；全部（轨迹，帧）观测合并的中位数与 95 分位。
\end{itemize}
\section{结果}\label{pD:sec:res}
\begin{table}[!htbp]\centering\small\caption{Sod：末帧密度误差（中位 / 95 分位）与覆盖率；不断轨时为 $<10^{-6}$ / 0.97\%。}
\resizebox{\linewidth}{!}{\begin{tabular}{lcccc}\toprule 情形 & 每粒子断轨次数 & 窗口化（旧） & 逐轨账本 & 出生轨迹 95 分位（逐轨）\\\midrule
随机断轨 1\%/帧 & 0.36 & $<10^{-6}$ / 3.1\%（95\%） & $<10^{-6}$ / 1.3\%（100\%） & 1.8\%\\
随机断轨 3\%/帧 & 1.06 & $<10^{-6}$ / 4.9\%（86\%） & $<10^{-6}$ / 1.6\%（100\%） & 2.1\%\\
随机断轨 5\%/帧 & 1.73 & 0.05\% / 6.6\%（79\%） & 0.01\% / 2.4\%（100\%） & 2.7\%\\
30\% 的粒子在穿越激波时断轨 & 0.46 & $<10^{-6}$ / 5.6\%（93\%） & $<10^{-6}$ / 1.9\%（100\%） & 3.9\%\\
1/4 稀疏 + 0.02$h$ 噪声 + 3\%/帧 & 1.11 & 2.7\% / 13.1\%（85\%） & 3.2\% / 12.8\%（100\%） & 12.7\%\\
\bottomrule\end{tabular}}\end{table}
\begin{table}[!htbp]\centering\small\caption{柱面爆炸波：末帧密度误差（中位 / 95 分位）与覆盖率；不断轨时为 0.58\% / 32.2\%。}
\resizebox{\linewidth}{!}{\begin{tabular}{lcccc}\toprule 情形 & 每粒子断轨次数 & 窗口化（旧） & 逐轨账本 & 出生轨迹 95 分位（逐轨）\\\midrule
随机断轨 1\%/帧 & 0.55 & 10.8\% / 68.1\%（95\%） & 1.1\% / 37.9\%（100\%） & 39.9\%\\
随机断轨 3\%/帧 & 1.61 & 12.7\% / 76.1\%（86\%） & 3.0\% / 48.0\%（100\%） & 49.5\%\\
随机断轨 5\%/帧 & 2.62 & 15.5\% / 93.4\%（78\%） & 6.5\% / 55.8\%（100\%） & 56.3\%\\
30\% 的粒子在穿越激波时断轨 & 0.74 & 12.0\% / 79.6\%（94\%） & 3.5\% / 43.4\%（100\%） & 44.9\%\\
1/4 稀疏 + 0.02$h$ 噪声 + 3\%/帧 & 1.63 & 23.5\% / 137.8\%（86\%） & 21.7\% / 127.9\%（100\%） & 131.5\%\\
\bottomrule\end{tabular}}\end{table}
\begin{table}[!htbp]\centering\small\caption{激波–等熵涡（DNS）：末帧密度误差（中位 / 95 分位）与覆盖率；不断轨时为 0.34\% / 3.4\%。}
\resizebox{\linewidth}{!}{\begin{tabular}{lcccc}\toprule 情形 & 每粒子断轨次数 & 窗口化（旧） & 逐轨账本 & 出生轨迹 95 分位（逐轨）\\\midrule
随机断轨 1\%/帧 & 0.27 & 0.39\% / 3.9\%（95\%） & 0.36\% / 3.6\%（100\%） & 2.9\%\\
随机断轨 3\%/帧 & 0.78 & 0.44\% / 4.2\%（86\%） & 0.38\% / 3.8\%（100\%） & 3.4\%\\
随机断轨 5\%/帧 & 1.27 & 0.48\% / 4.5\%（78\%） & 0.40\% / 4.0\%（100\%） & 3.6\%\\
30\% 的粒子在穿越激波时断轨 & 0.56 & 0.42\% / 4.3\%（91\%） & 0.40\% / 4.1\%（100\%） & 3.3\%\\
1/4 稀疏 + 0.02$h$ 噪声 + 3\%/帧 & 0.80 & 2.9\% / 12.7\%（86\%） & 3.5\% / 14.3\%（100\%） & 13.6\%\\
\bottomrule\end{tabular}}\end{table}
\begin{figure}[!htbp]\centering\includegraphics[width=\linewidth]{figS.pdf}\caption{末帧密度误差的 95 分位。灰柱上方的数字为窗口化方法的覆盖率（逐轨方法均为 100\%）；虚线为不断轨时的账本。}\end{figure}
\paragraph{Sod。}随机断轨 1\%、3\%、5\%/帧时，逐轨账本（延迟标定）的末帧 95 分位为 1.3\%、1.6\%、2.4\%，窗口化为 3.1\%、4.9\%、6.6\%，后者的覆盖率只有 95\%、86\%、79\%。激波处断轨的对抗情形中，逐轨为 1.9\%，窗口化为 5.6\%（覆盖率 93\%）。
\paragraph{柱面爆炸波。}在 5 个情形中的 5 个里，逐轨账本（延迟标定）的 95 分位不高于窗口化；逐轨覆盖率均为 100\%，窗口化为 78\%–95\%。激波处断轨情形：逐轨 43.4\%，窗口化 79.6\%。
\paragraph{激波–等熵涡（DNS）。}在 5 个情形中的 4 个里，逐轨账本（延迟标定）的 95 分位不高于窗口化；逐轨覆盖率均为 100\%，窗口化为 78\%–95\%。激波处断轨情形：逐轨 4.1\%，窗口化 4.3\%。
\paragraph{同一批轨迹上的公平比较（激波–涡，3\%/帧）。}窗口化只对窗口内完整的轨迹给出估计，所以上表的 95 分位不是在同一批轨迹上比较的。只取两种方法都有估计的轨迹（10412 条）：逐轨账本中位 0.42\%、95 分位 4.2\%，窗口化中位 0.44\%、95 分位 4.2\%，两者相当；逐轨账本另外覆盖了窗口化漏掉的 1673 条轨迹，其 95 分位为 1.9\%。涡核内（$r_0<2R_v$）两者的 95 分位为 7.7\% 与 7.1\%，这是账本在涡核内的固有误差（见噪声与锚定的分析），与断轨无关。
\paragraph{爆炸波是最难的情形。}不断轨时账本的 95 分位已达 32.2\%（激波正扫过的薄壳）。断轨后，逐轨账本的中位误差为 1.1\%–6.5\%，窗口化为 10.8\%–15.5\%；95 分位分别为 37.9\%–55.8\% 与 68.1\%–93.4\%。逐轨账本好得多，但 95 分位仍高：61 帧内每个粒子平均断轨 0.6–2.6 次，中心区膨胀使示踪点变稀，而新出现的轨迹常落在密度为 6 的薄壳内，用近邻标定时会跨越陡峭的密度梯度。稀疏加噪声时两种方法都失效（95 分位超过 100\%）：稀疏后的间距已与薄壳厚度相当。
\paragraph{开发中修掉的第二个问题。}第二版在激波外遇到不可测的环时沿用上一帧的值（丢掉这一步的增量），并在涡核内长期使用缩减环（缩减后的多边形更大，非仿射误差更大）。改为“所有断环都在锚定帧重建、激波外丢顶点后立即重锚”后，激波–涡在同一批轨迹上的 95 分位从 6.1\%% 降到 4.2\%%。\n
\section{结论与局限}
\begin{enumerate}
\item 逐轨账本让每条轨迹在其连续段内保持自己的锚定星：顶点消失时缩减环，激波中的断环在锚定帧重建，锚定点因此不会被断轨拖进激波。
\item 覆盖率由窗口化的 79\%–95\% 提高到 100\%：末帧存在的每条轨迹都有密度估计，包括断轨后重新出现的轨迹。
\item 出生轨迹的标定是主要的剩余误差来源；延迟到离开激波模板之后再标定，总体上比即时标定好。
\item 稀疏加噪声的组合情形中，误差由噪声主导（见噪声传播律），逐轨与窗口化差别不大，但逐轨的覆盖率仍为 100\%。
\item 爆炸波的薄壳仍是难点：新出现的轨迹在陡峭密度梯度中用近邻标定，误差大；改进方向是用账本的质量守恒（局部质量平衡）代替近邻插值来标定新轨迹。
\item 仍需要第 0 帧（或某一参考时刻）的密度；整条轨迹全程都不与已标定轨迹相邻的孤立出生轨迹无法标定。
\end{enumerate}

\clearpage\part{第五部分　场图可视化：输运、平流、扩散与掺混}
\noindent\fcolorbox{accent!40}{tint}{\begin{minipage}{0.965\linewidth}\vspace{0.3em}
\textbf{\sffamily 本部分的内容与读法}\quad 每个算例一节，每节先给出“读图”（逐行说明每一行画的是什么、颜色代表什么），再给出“定量读数”（从图中各场直接算出的数，脚本 \texttt{scripts/viz\_stats.py}，结果 \texttt{results/res\_viz\_stats.json}），最后是“物理解释”与“伪影与局限”。图中的场全部只用示踪点位置（加初始密度）计算，标注为 DNS 的行仅作对照。通用约定如下：
\begin{itemize}[leftmargin=1.4em,itemsep=0.1em]
\item \textbf{物质标签图}：按初始归属给示踪点着色（红、蓝两类物质），黑点为初始位于某条物质线或界面上的示踪点。它显示的是\emph{平流与物质输运}本身：红色物质跑到哪里，就有多少物质被输运过去。
\item \textbf{扩散标量图}：初始为一条光滑锋面的被动标量，在网络上按质量形式的拉格朗日扩散方程演化（T8）。锋面被拉长、变薄并展宽，显示的是\emph{平流与扩散的共同作用}，即混合。
\item \textbf{拉伸场 $\ln\lambda$}：链式张量估计的初始标签等值线族的切向拉伸，红为拉长、蓝为缩短；\textbf{压缩场 $-\ln J$}：账本给出的面积压缩，橙为压缩、紫为膨胀。二者之和就是界面上标量梯度的增长（Nanson 恒等式），即\emph{掺混的两个来源}。
\item \textbf{各向异性 $A=\sigma_{\max}(F)/\lambda$}：大值表示标签等值线不在最大拉伸方向上，链式张量在这里对误差敏感。
\end{itemize}
\vspace{0.2em}\end{minipage}}

\section{双环流（不可压，解析速度场）}
\paragraph{读图。}第一行：左环流的物质为红、右环流为蓝，黑点是初始物质线 $x_0=0.5$；第二行：初始锋面在 $x=1$ 的扩散标量（$\kappa=10^{-3}$）；第三行：只凭位置的步涡量 $\hat\omega=2A\,\mathrm{Im}\,p/(\bar A\Delta t)$；第四行：$t=5$ 时由位置差分得到的速度矢量，以及 $t=10$ 时的拉伸场、各向异性与压缩场。四列为 $t=0,3.4,6.6,10$。
\paragraph{定量读数。}左环流的物质进入右半区的比例依次为 0、14.6\%、20.0\%、32.9\%：分界线 $x\approx1$ 的周期摆动每次把一个“瓣”的物质卷入对侧。扩散标量处于混合区（$0.1<c<0.9$）的示踪点比例为 11.3\%、18.5\%、35.1\%、59.2\%，按质量加权的方差比为 1、0.93、0.77、0.56。步涡量在 $t\approx6.6$ 的 1\%–99\% 分位为 $-2.5$ 到 $1.7$，两个环流的符号相反。$t=10$ 时拉伸场的中位数为 0.96，5\% 与 95\% 分位为 $-0.74$ 与 3.14，22.6\% 的示踪点 $|\ln\lambda|>2$；跟踪的物质线上平均拉伸为 1.708（逐点精确值 1.711）。压缩场的 95\% 分位只有 0.06。
\paragraph{物理解释。}这是典型的瓣输运与 Batchelor 型混合：物质每个周期被折叠一次，物质线被拉长、折叠成 S 形，扩散标量的锋面沿物质丝被拉薄，梯度变陡，扩散随之加速，所以方差在后半段衰减得更快（$t=6.6\to10$ 的方差比下降 0.21，前 3.4 个时间单位只下降 0.07）。各向异性的中位数为 1.36、95\% 分位为 14.4，且 $|\ln\lambda|$ 与 $A$ 的秩相关为 $-0.45$：拉伸最强的地方，等值线恰好沿着最大拉伸方向（$A\approx1$）；$A$ 大的地方是等值线被侧向剪切、切向反而不长的区域。
\paragraph{伪影与局限。}不可压流的压缩场应恒为零，图中 $\pm0.06$ 的斑点是账本在强折叠物质丝上的离散误差，与 T2 的 95 分位（5.7\%）一致。
\begin{figure}[p]\centering\includegraphics[width=\linewidth]{viz_dg.pdf}
\caption{双环流的场图。第一行：物质标签与物质线 $x_0=0.5$；第二行：网络扩散标量（质量形式）；第三行：只凭位置的步涡量；第四行：位置差分速度、拉伸场、各向异性与压缩场。}\end{figure}

\clearpage\section{Lamb–Oseen 共转涡对（不可压）}
\paragraph{读图。}第一行：$x<0$ 的物质为红，黑点是绕左涡、半径 1.5 的物质圆环；第二行：初始锋面在 $x=0$ 的扩散标量（$\kappa=0.02$）；第三行：步涡量；第四行：速度矢量与末时刻（$t=6\pi$）以绕左涡圆周为切向的拉伸场、各向异性和压缩场。四列为 $t=0,6.1,12.7,18.8$。
\paragraph{定量读数。}原在左半区的物质进入右半区的比例为 0、15.9\%、31.1\%、40.5\%，这主要是涡对整体公转带来的输运，而不是混合；混合区比例为 6.5\%、14.8\%、22.2\%、31.7\%，方差比为 1、0.93、0.85、0.76，混合明显慢于双环流。步涡量的峰值为 1.81（解析峰值 $\Gamma/\pi a^2=2$，三角形平均使峰值略低）。物质圆环到末时刻仍近圆：相对其中心的半径的相对标准差为 4.6\%；圆环上的局部拉伸在 $-2.31$ 到 $0.55$ 之间（5\%–95\% 分位），平均 $-0.59$。圆环上各向异性的中位数为 14.1，全场只有 2.0。
\paragraph{物理解释。}涡对的输运以整体旋转为主，两臂螺旋把物质卷绕进去，但涡核本身保持相干，所以扩散标量的锋面只是被卷绕、没有被大量拉薄，方差衰减慢。物质圆环在差动旋转下总长几乎不变，但局部被另一个涡的应变时而拉长、时而缩短，这就是第一部分“圆环拉伸约为 0”被推翻的直观原因：平均为负，局部在 $-2.3$ 到 $0.6$ 之间。
\paragraph{伪影与局限。}圆环恰处在强剪切中（$A\approx14$），链式张量的误差在这里被放大（第一部分：$N=8000$ 时界面平均 $-0.55$，精确 $-0.31$，随分辨率二阶收敛）；压缩场 95\% 分位 0.13 也来自涡周星的非仿射变形，自适应锚定把这一区域的密度 95 分位从 13\% 降到 7\%。
\begin{figure}[p]\centering\includegraphics[width=\linewidth]{viz_lo.pdf}
\caption{Lamb–Oseen 共转涡对的场图。各行同双环流；第四行的拉伸场以绕左涡的圆周为切向。}\end{figure}

\clearpage\section{Sod 激波管（可压，精确轨迹）}
\paragraph{读图。}第一行：驱动气体（$x_0<0$，红）与被驱动气体（蓝），黑点为物质线 $x_0=0.1$；第二行：只凭位置的账本密度；第三行：账本密度的相对误差（色标 $\pm5\%$）；第四行：初始锋面在接触面的扩散标量（$\kappa=10^{-3}$）；第五行：C 通道检出的穿越时刻（画在初始位置上）、$t=0.14$ 的步膨胀率 $\hat\theta$、同一时刻的步涡量（理论上应为零）与速度矢量。四列为 $t=0.04,0.105,0.175,0.24$。
\paragraph{定量读数。}由检出的穿越时刻对初始位置线性拟合，得激波速度 1.755（精确 1.752）；被检出穿越的示踪点占 35\%。账本给出的波后平台密度为 0.26557，接触面左侧为 0.42632，与精确解的 0.26557、0.42632 一致到 5 位；相对误差超过 1\% 的示踪点只占 5.1\%，全部在激波与接触间断两侧的模板内。扩散锋面的 10\%–90\% 宽度依次为 0.061、0.071、0.077、0.080。激波处步膨胀率达到 $-151$，同一步的跨激波三角形上伪涡量的最大幅值为 23，约为膨胀率的 15\%。
\paragraph{物理解释。}物质标签图直接显示接触面就是两种气体的物质边界，以速度 0.93 右移；物质线 $x_0=0.1$ 在激波扫过时被一次性压缩（压缩 0.754），之后随流平移，拉伸恒为零——平面激波对掺混的贡献完全来自压缩。扩散锋面的宽度增长随时间变慢，是一维扩散 $\sqrt{\kappa t}$ 的表现；锋面在接触面两侧的密度不同（0.43 与 0.27），质量形式正确计入了这一密度差（体积形式误差大一个数量级，见第二部分）。
\paragraph{伪影与局限。}第五行第三个子图里沿激波的条带是伪涡量：跨激波三角形的三点映射不是单轴的，$\mathrm{Im}\,p\neq0$。这是“跨激波只报告积分量、用合并块计算转角”的直观理由；合并块后伪转角的 95 分位从 0.092 rad 降到 0.0057 rad。
\begin{figure}[p]\centering\includegraphics[width=\linewidth]{viz_sod.pdf}
\caption{Sod 激波管的场图。第一至四行：物质标签、账本密度、密度误差、扩散标量；第五行：穿越时刻、步膨胀率、伪涡量与速度。}\end{figure}

\clearpage\section{柱面爆炸波（可压，精确轨迹）}
\paragraph{读图。}第一行：$r_0\le0.2$ 的高密度内部气体为红，黑点为物质圆环 $r_0=0.5$；第二行：账本密度（色标 0–6）；第三行：初始锋面在 $r=0.5$ 的扩散标量（$\kappa=2\times10^{-4}$）；第四行：穿越时刻、$t=25$ 的拉伸场（切向取圆周方向）与压缩场、$t=13$ 的速度矢量。四列为 $t=1,9,17,25$。
\paragraph{定量读数。}对检出的穿越时刻与初始半径做双对数拟合，指数为 0.54（理论的 $r_s\propto t^{1/2}$ 给出 0.5）。末时刻外部气体被压进 $r\approx0.94$–$1.00$ 的薄壳，账本密度的 99.5\% 分位为 5.73（壳内精确峰值约 6）；壳内压缩的中位数 1.48。内部气体强烈膨胀：中心区 $-\ln J$ 为 $-2.68$，即密度降到初始的 $1/14.6$，内部气体的中位半径从 0.2 增大到 0.54。圆环上拉伸 0.576（精确 0.578）、压缩 0.637（精确 0.638）。扩散标量的混合区比例为 6.1\%、21.0\%、43.9\%、55.5\%，方差比为 1、0.85、0.65、0.56。
\paragraph{物理解释。}扩散标量的环形锋面先被激波压进薄壳而变陡（$t=9$ 混合区已达 21\%），随后在膨胀流中被拉长并展宽；方差在激波扫过锋面的阶段衰减最快。圆环同时受到两种作用：激波一次性的压缩（使梯度变陡）与之后随流膨胀的径向拉伸（使圆周变长），所以界面梯度的增长是“先压缩、后被膨胀部分抵消”。拉伸场在全场为径向自相似的均匀拉伸，没有局部结构，与流动的一维性质一致。
\paragraph{伪影与局限。}薄壳只有约三个示踪点间距厚，激波正扫过的示踪点的星跨在激波上，这是账本 95 分位（32\%）的来源；第四部分的断轨检验中，爆炸波的薄壳也是唯一仍有大误差的区域。
\begin{figure}[p]\centering\includegraphics[width=\linewidth]{viz_blast.pdf}
\caption{柱面爆炸波的场图。第一至三行：物质标签与圆环、账本密度、扩散标量；第四行：穿越时刻、拉伸场、压缩场与速度。}\end{figure}

\clearpage\section{激波–等熵涡（可压，DNS）}
\paragraph{读图。}第一行：初始涡核（$r_0<R_v$）的物质为红，黑点为 $r_0=2R_v$ 的物质圆环；第二行：账本密度；第三行：账本密度相对 DNS 的误差（色标 $\pm5\%$）；第四行：只凭位置的步涡量；第五行：DNS 涡量（对照）。四列为 $t=0,0.2,0.4,0.6$。
\paragraph{定量读数。}涡核物质的质心位移为 0、0.0004、0.060、0.149：激波在 $t\approx0.2$ 到达涡核之前涡核几乎不动，之后被波后流动带走。圆环的平均压缩为 0、0.12、0.42、0.40，激波扫过后基本不变。账本密度相对 DNS 的中位误差为 0.09\%、0.22\%、0.32\%，95 分位为 1.1\%、2.4\%、3.2\%。$t=0.4$ 时误差最大的 5\% 中有 27\% 位于近期的激波模板内，末时刻有 93\% 位于涡核内。步涡量与 DNS 涡量在示踪点上的相关系数为 0.994、0.998、0.999（三个时刻）。
\paragraph{物理解释。}激波扫过等熵涡时，涡核物质保持相干、被压缩成椭圆并随波后流动输运；涡量几乎不被激波改变（等熵涡没有密度梯度与压力梯度的交叉，斜压沉积很小），所以位置涡量与 DNS 涡量在激波前后都高度一致。误差的位置随时间迁移：激波经过时集中在激波带，之后集中在涡核——后者是涡核内锚定星的非仿射变形，第三部分的 DNS 加密已确认它与 DNS 网格无关。
\paragraph{伪影与局限。}激波正扫过涡核的帧上，第四行可以看到沿激波的伪涡量条带，第五行的 DNS 涡量里没有；在这些三角形上应只报告积分量。
\begin{figure}[p]\centering\includegraphics[width=\linewidth]{viz_sv.pdf}
\caption{激波–等熵涡的场图。第一至三行：涡核物质与圆环、账本密度、密度误差；第四、五行：只凭位置的步涡量与 DNS 涡量。}\end{figure}

\clearpage\section{激波–重气泡（可压，无粘 DNS）}
\paragraph{读图。}第一行：重气体示踪点（红）与界面示踪点（黑）；第二、三行：账本密度与 DNS 密度；第四、五行：网络扩散标量（质量形式）与带同一扩散系数的 DNS 标量；第六行：$t=1.4$ 的拉伸场（切向取初始圆周方向）、压缩场、各向异性与 $t=1.38$ 的步涡量。四列为 $t=0,0.6,1.0,1.4$。
\paragraph{定量读数。}重气体的流向范围从 $[0.52,0.98]$ 变为 $[0.71,0.98]$（$t=0.6$，被压扁）、$[0.87,1.15]$、$[1.07,1.42]$（$t=1.4$，卷成肾形并整体下游输运）；最高点从 0.24 先降到 0.22，再被卷起抬到 0.30。账本密度相对 DNS 的中位误差为 0.28\%、0.58\%、0.67\%，末时刻误差最大的 5\% 中有 87\% 位于距界面 0.03 以内。网络扩散标量与 DNS 标量的相对 $L^2$ 差为 1.4\%、2.3\%、7.3\%；混合区比例几乎相同（网络 23.7\%、33.1\%、40.2\%，DNS 24.0\%、32.9\%、40.4\%）。$t=1.4$ 时界面上的拉伸在 $-1.50$ 到 1.53 之间（5\%–95\% 分位），平均 0.29；压缩在界面上平均 0.47，在重气体内部 0.63，在轻气体中 0.58。界面上各向异性的中位数为 5.0，远离界面为 2.1。步涡量的 1\%–99\% 分位为 $-18.2$ 到 3.3，以负值（上半平面顺时针）为主。
\paragraph{物理解释。}物质标签图给出了输运的完整图像：激波先把气泡沿流向压扁（$t=0.6$），斜压沉积的涡片随后把重气体卷成涡对（$t=1.0$–1.4），重气体沿螺旋臂进入涡中。拉伸场与压缩场把掺混分成两部分：压缩是激波给的、几乎均匀（重气体内 0.63）；拉伸集中在螺旋臂上，有正有负——卷起的臂被拉长（红），与激波大致平行的界面段被切向缩短（蓝）。这正是第三部分“激波扫过后界面先被切向缩短、抵消大部分压缩”的空间图像。步涡量以负值为主，对应上半气泡顺时针的斜压环量（第三部分：位置形式的上半气泡环量 0.266，与 DNS 相差 1.5\%，与 Samtaney–Zabusky 模型相差 8\%）。网络扩散标量与 DNS 标量的混合区几乎一致，说明质量形式的网络扩散在卷起的螺旋中仍然可靠。
\paragraph{伪影与局限。}界面附近的密度误差主要来自 DNS 本身的数值扩散（第二部分的 DNS 加密）；卷起后的涡片上各向异性大（$A\approx5$），链式张量在这里误差被放大，且无粘 DNS 的界面带不收敛（信息极限）。
\begin{figure}[p]\centering\includegraphics[width=\linewidth]{viz_sb.pdf}
\caption{激波–重气泡（无粘 DNS）的场图。第一至五行：重气体与界面、账本密度、DNS 密度、网络扩散标量、DNS 扩散标量；第六行：拉伸场、压缩场、各向异性与步涡量。}\end{figure}

\clearpage\section{粘性激波–重气泡（$Re\approx1300$）}
\paragraph{读图。}第一行：重气体与界面，$t=0.6,1.0,1.4$；第二行：账本密度；第三行：$t=1.4$ 的密度相对误差（对粘性 DNS，色标 $\pm10\%$）、拉伸场与压缩场。
\paragraph{定量读数。}重气体的整体输运与无粘算例相近（$t=1.4$ 流向范围 $[1.08,1.40]$，最高点 0.297），但界面光滑、没有细小卷起。账本密度相对粘性 DNS 的中位误差为 0.27\%、0.43\%、0.48\%，误差最大的 5\% 中有 79\% 位于界面附近。界面上的拉伸在 $-1.12$ 到 1.30 之间，平均 $-0.02$；压缩在界面上平均 0.32，在重、轻气体内部为 0.63 与 0.58；界面各向异性的中位数为 3.7（无粘为 5.0）。
\paragraph{物理解释。}物理粘性抑制了涡片的卷起，界面整体没有被拉长（平均拉伸接近零），所以到 $t=1.4$ 界面上梯度的增长几乎全部来自压缩。与无粘算例对比，两者的整体输运相同，差别在界面的小尺度结构：这是“压缩与拉伸的分配依赖雷诺数”的直观证据。界面压缩（0.32）低于两侧内部（约 0.6），是因为物理扩散在界面处把两种气体的密度混合，界面上的物质元在压缩后又随混合膨胀。
\begin{figure}[p]\centering\includegraphics[width=\linewidth]{viz_sbv.pdf}
\caption{粘性激波–重气泡（$\Delta x=0.0035$）的场图。第一、二行：重气体与界面、账本密度；第三行：密度误差、拉伸场与压缩场。}\end{figure}

\clearpage\section{单模 Richtmyer–Meshkov（可压，DNS）}
\paragraph{读图。}第一行：重气体（$x>0.5+a_0\cos2\pi y$，红）与界面示踪点；第二行：账本密度；第三行：步涡量；第四行：$t=2$ 的密度误差、拉伸场（切向取初始界面方向）、压缩场与 $t=0.4$ 的速度矢量。四列为 $t=0,0.6,1.2,2.0$。
\paragraph{定量读数。}激波扫过后界面以 $\Delta u\approx0.51$（速度矢量的中位数）整体右移，扰动振幅先被压缩、随后线性增长（第二部分：增长率 0.042，Richtmyer 模型 0.043）。$t\approx1.2$ 时距界面 0.03 以内的示踪点承载了 65\% 的 $|\hat\omega|$ 积分：涡量集中在界面上的斜压涡片中。$t=2$ 时界面压缩平均 0.80，重、轻气体内部为 0.76 与 0.78；拉伸在尖钉处（$y\approx0$）为 $+0.34$，气泡顶（$y\approx0.5$）为 $+0.07$，侧翼（$y\approx0.25$）为 $-0.12$。远离界面的密度中位误差为 0.44\%。
\paragraph{物理解释。}重气体的尖钉在 $y=0$ 处伸入轻气体、轻气体的气泡在 $y=0.5$ 处伸入重气体，界面上的涡片驱动这一增长。拉伸场显示界面被拉长的位置是尖钉与气泡顶，而侧翼被缩短；但到 $ka\approx0.6$ 拉伸的幅度（0.1–0.3）仍远小于激波的压缩（0.8），所以这一阶段界面梯度的增长约 94\% 来自压缩。压缩场在两侧气体中几乎相同，说明激波压缩是整体性的。
\paragraph{伪影与局限。}紧贴界面（$|\phi_0|<0.02$）的密度“误差”中位数高达 35\%，但这不是账本的误差：DNS 在接触间断上有数值扩散，把两种气体的密度混在一起（与激波–重气泡的 DNS 加密结论相同），界面上的 DNS 密度不能作参照；远离界面账本与 DNS 一致到 0.44\%。
\begin{figure}[p]\centering\includegraphics[width=\linewidth]{viz_rm.pdf}
\caption{单模 Richtmyer–Meshkov 的场图。第一至三行：重气体与界面、账本密度、步涡量；第四行：密度误差、拉伸场、压缩场与速度。}\end{figure}

\clearpage\section{小结：图中能看到的五个机理}
\begin{enumerate}
\item \textbf{平流与物质输运}：物质标签图把输运直接画了出来（瓣输运、涡对卷绕、激波后整体平移、气泡卷成涡对、RM 的尖钉与气泡），并可按质量计量（双环流 32.9\%、Lamb–Oseen 40.5\% 的物质跨过中线）。
\item \textbf{压缩}：激波在物质上留下一次性的、几乎均匀的压缩（Sod 0.754、激波–重气泡约 0.6、RM 约 0.8），膨胀则表现为负的 $-\ln J$（爆炸波中心 $-2.7$）。
\item \textbf{拉伸}：只在有涡或剪切的地方出现，且有正有负——双环流的物质丝、Lamb–Oseen 的涡周、激波–重气泡的螺旋臂、RM 的尖钉。
\item \textbf{扩散与混合}：扩散标量在平流拉薄的地方混合最快（双环流后半段、爆炸波激波扫过锋面时），网络质量形式与 DNS 标量的混合区几乎一致。
\item \textbf{方法的边界}：跨激波三角形的伪涡量条带、涡周与涡片上的高各向异性、接触间断上 DNS 的数值扩散，都在图中直接可见。
\end{enumerate}


\clearpage\part{第六部分　v2：修正、补证与验证}
\noindent\fcolorbox{accent!40}{tint}{\begin{minipage}{0.965\linewidth}\vspace{0.4em}
\textbf{\sffamily 本部分的结论}\quad
(1) 九条代数恒等式全部通过精确符号验证；定理 1 推广到任意多边形，单轴签名证明为充要条件，定理 2 改称运动学恒等式并补出带误差界的离散 Kelvin 推论。
(2) 在制造流映射上：平面激波处账本精确到 $10^{-13}$ 量级；弯曲激波处误差只取决于 $h/R_s$；光滑流中多边形账本是一阶，新增的点态张量账本是二阶。
(3) 同区域环量比较：在同一物质边界上，位置形式环量与 DNS 速度线积分的中位相差为 $0.03\%$–$0.25\%$，四个 DNS 算例都如此；旧的“与 DNS 相差 1.5\%”比较的是不同区域，已废弃。
(4) 扩散的“质量形式准 4–10 倍”实为 3.8--10.6 倍，差距来自方程本身，不来自账本；在 $\rho D\propto T^{0.7}$ 下差距依然存在。
(5) 账本 v3（多窗口激波检测、自动推迟上限、张量/混合估计）：弱激波和慢激波由失效变为精确；除 Sod（中位误差本已是机器精度）外，六个流动算例的中位误差下降 36\%–91\%；激波处的 95 分位基本不变。
\vspace{0.3em}\end{minipage}}

\section{本轮修正的不一致}
\begin{table}[!htbp]\centering\small\caption{论文与报告中已修正的表述。数字由 \texttt{tools/paper\_numbers.py} 从结果文件计算，\texttt{tools/check\_paper\_numbers.py} 核对：v1 有 34 处问题，v2 全部通过。}
\begin{tabularx}{\linewidth}{@{}l L L@{}}\toprule
位置 & 旧表述 & 现行表述\\\midrule
总摘要 & 密度中位误差 $<1\%$ & 至多 $1.3\%$（双环流 1.3\%，Sod 为机器精度）\\
扩散 & 质量形式比体积形式准 4–10 倍 & 3.8--10.6 倍（Sod 10.6、爆炸波 10.2、激波–气泡 3.8）；误差定义为 RMS 误差除以参照的标准差\\
环量 & 与 DNS 相差 1.5\%（物质带对整个上半区域） & 与同一物质边界上的 DNS 速度线积分比较，见表~\ref{v2:tab:circ}\\
Kelvin & 物质带环量守恒 & 只有宽带近似守恒；窄带在位置形式和线积分中一致下降（0.257$\to$0.159），原因是无粘 DNS 在涡片附近的数值耗散\\
参数敏感性 & 锚定周期 3–8 帧对中位误差影响 $<0.03$ 个百分点 & 锚定周期取 3 帧时，爆炸波中位误差 0.58\%$\to$0.77\%、95 分位 32\%$\to$46\%；$\tau$ 的影响 $\le$0.013 个百分点\\
RM & 增长率 0.0421 & 0.0423（由振幅序列在 $t\in[0.6,1.4]$ 上拟合）；Richtmyer 参数与一维激波折射精确解（$\Delta u=0.5112$，$\rho=2.178/6.413$，$A^+=0.493$）相差 $\le0.7\%$\\
T9 估计量 & 两者同时失效“当且仅当” & 改为充分条件，并给出信息极限的数值（第~\ref{v2:sec:mfm} 节）\\\bottomrule
\end{tabularx}\end{table}

\section{证明的补强}
\begin{theorem}[有限步环量与面积变化，多边形形式]\status{已证+符号验证}
记 $u^{\rm m}_e=\frac12(\Delta x_a+\Delta x_b)$，$q^{\rm m}_e=\frac12(q^n_e+q^{n+1}_e)$，$c_e=\overline{u^{\rm m}_e}\,q^{\rm m}_e$（上横线只表示共轭）。对任意由示踪点构成的闭合多边形 $\Pi$，有 $\sum_{\partial\Pi}\im c_e=A^{n+1}_\Pi-A^n_\Pi$；且对 $\Pi$ 的任意三角剖分，$\sum_{\partial\Pi}\re c_e=\sum_t2A^n_t\im p_t$。
\end{theorem}
证明的要点：顶点为 $x^n+s\Delta x$ 的多边形，其鞋带面积是 $s$ 的二次式，所以 $A(1)-A(0)=A'(\tfrac12)$ 精确成立。虚部就是 ALE 方法中的二维离散几何守恒律（Thomas–Lombard 1979；Farhat 等 2001），本工作的新意在于把它与环量配成一个复数，并用作质量账本。

\begin{corollary}[离散 Kelvin 定理]\status{已证}
定理 4.2（位置形式）本身是时间上的分部求和，不含动力学，所以改称“离散运动学环量恒等式”。若示踪点沿光滑流映射运动，且 $a=-\nabla\Psi$（正压、无粘），则对周长为 $L$、最长边为 $h$ 的物质多边形，
$\big|\sum_{t\in R}(\Gamma^{n+1/2}_t-\Gamma^{n-1/2}_t)\big|\le\frac{\Delta t^2L}{12}\big(h^2\|\nabla^2a\|_\infty+\Delta t^2\|D_t^2a\|_\infty\big)+O(\Delta t^6L)$。
\end{corollary}
推论的证明：主项是 $\oint a\cdot dx=0$ 的梯形公式，每条边的误差不超过 $\ell^3\|\nabla^2a\|/12$。

\begin{proposition}[单轴签名是充要条件]\status{已证+符号验证}
对迹为正的 $F$，$\arg p=0$ 当且仅当 $F$ 对称；对对称正定的 $F$，$|\mu|=\tanh(\frac12|\ln J|)$ 当且仅当某个特征值等于 1。两者之差的平方的分子分解为 $-4\lambda_1\lambda_2(\lambda_1-1)(\lambda_1+1)(\lambda_2-1)(\lambda_2+1)$。
\end{proposition}
另有两点。第一，合并块就等于块两端帧之间的直接三点映射。第二，偏离签名量度的是激波曲率和波后流动的不均匀性（弯曲激波上的涡量跃变，Hayes 1957）；在会聚柱面激波上，$|\arg p_B|\approx0.16\,h/R_s$。

\paragraph{账本估计的是什么。}账本给出 $\hat\rho_i=\rho_i(t_c)A^\ast_i(t_c)/A^\ast_i(t)$，目标是 $\rho(X_i,t_c)/J(X_i,t)$。多边形面积是顶点的二次型，所以对任意仿射预测 $\hat x$ 和残差 $e$，精确成立
$A^\ast(\hat x+e)=\det F\,A^\ast(t_a)+\sum_k\frac12e_k\times(\hat x_{k+1}-\hat x_{k-1})+A^\ast(e)$。
由此可得四点：
\begin{itemize}
\item 星作仿射运动时，账本精确。
\item 取 $\delta=\ln[A^\ast(t)/A^\ast(t_a)]-\ln\det F_{\rm LS}$，它恰好是一个增量中仿射运动解释不了的部分，可作为逐点指示量。
\item 噪声律是同一展开取 $e$ 为噪声的特例。
\item 光滑流中，多边形面积比估计的是星面积质心处的 $J$，而不是示踪点处的 $J$，所以账本是一阶。
\end{itemize}
望远相消命题的假设写精确为：平面激波速度恒定，两侧为均匀态。在这些假设下，穿越时刻是初始位置的仿射函数，所以无论激波用了几帧扫过星，面积比都精确。

\paragraph{Dirac 一节的表述。}定理 1、2 与账本只用到与度量无关的 $d_1$，度量只进入 $L_0$ 与 T7；$M_1$ 取 Whitney 质量阵（对角余切 Hodge 星给出同一个 $d_0^\top M_1d_0$，但 $L_1$ 不同）。方差无条件衰减的前提补为 $\rho\kappa\ge0$ 且三角形正定向。

\paragraph{M 矩阵版本。}若改用边电导（相邻三角形电导的调和平均），并把剥离边界边上为负的余切权截断为零，$K$ 就是 M 矩阵，对任何电导场都满足离散极值原理且半正定。注意不截断时它不是半正定的，实测会发散。实际效果是：单元形式的越界量 $\le3\times10^{-8}$，而截断版在狭长的 Sod 区域上误差升到 1.77\%（单元形式为 0.17\%），所以默认仍用单元形式。

\paragraph{T9 的补充。}$\kappa>0$ 时等值线不再是物质线，Nanson 分解只是梯度收支的运动学部分。薄层的标量耗散 $\propto\rho\kappa\lambda^2/J$，因此按积分混合率计算时拉伸项要计两次；RM 中为 $2\times0.054$ 对 $0.80$，压缩仍约占 88\%。

\section{制造流映射上的验证}\label{v2:sec:mfm}
为把方法误差与 DNS 误差分开，我们直接指定解析流映射，使密度、$J$ 与拉伸都有精确值（\texttt{lhd/mfm.py}）。结果见图~\ref{v2:fig:verif}。
\begin{itemize}
\item \textbf{平面激波}（$\theta\in\{0,0.5\}$，$\eta\in\{1.5,3,6\}$，含切向滑移，$C_s=0.5$–$4$）：只要激波被检出，账本误差 $\le6\times10^{-14}$，合并块单轴到 $7\times10^{-13}$；不推迟重锚时误差最大 41\%。单步检测会漏掉弱激波和慢激波：$\eta=1.5$、$C_s=0.5$ 时最大误差 12\%；$\eta=3$、$C_s=0.25$ 时 58\%。
\item \textbf{会聚柱面激波}（后方为汇流）：三种 $h$ 的误差曲线按 $h/R_s$ 重合，中位误差 0.08–0.34\%（$h/R_s=0.01$–$0.12$）。
\item \textbf{光滑可压缩旋转映射}：多边形账本观测阶 1.10，张量账本观测阶 2.18。
\item \textbf{薄剪切层}：$\delta/h=0.5,1,2,4,8$ 时，链式张量拉伸的相对误差为 11.9\%、10.9\%、5.0\%、2.0\%、0.7\%，信息极限约为 $\delta\gtrsim2$–$4h$。
\item \textbf{变形关系检验}（Sod）：刚体运动（含随时间变化的转动）、编号置换与缩放下，$\ln J$ 的最大变化为 1.5$\times10^{-12}$；时间反演恢复初始密度的中位误差为机器精度，95 分位 1.1\%。
\end{itemize}
\begin{figure}[!htbp]\centering\includegraphics[width=\linewidth]{fig_verification.pdf}
\caption{验证（图中文字为英文，与论文 v2 相同）。(a) 平面制造激波的最大密度误差；(b) 会聚柱面激波：账本中位误差与块转角随 $h/R_s$ 的变化；(c) 光滑映射：多边形账本与张量账本；(d) 激波–气泡环量：位置形式、同一边界上的 DNS 线积分、旧的顶点涡量求积；(e) 扩散：电导中用精确、账本、SPH、DTFE 密度，以及面积度量；(f) 薄剪切层的信息极限。}\label{v2:fig:verif}\end{figure}

\section{同区域环量比较}
旧比较把物质带的位置形式环量与 DNS 在另一区域上的涡量积分相比。正确的参照是 DNS 速度沿同一条物质多边形边界的线积分。为此重跑了四个 DNS（激波–等熵涡补存了示踪点速度）。
\begin{table}[!htbp]\centering\small\caption{位置形式环量对同边界 DNS 线积分的相对差（取 $|\Gamma|$ 超过其最大值 20\% 的时刻）；最后一列是旧参照（顶点涡量面积求积）对线积分的相对差。}\label{v2:tab:circ}
\begin{tabularx}{\linewidth}{@{}L r r r@{}}\toprule
算例与区域 & 中位（\%） & 最大（\%） & 旧参照中位（\%）\\\midrule
无粘激波–气泡，$|\phi_0|<0.15$ & 0.04 & 0.9 & 3.3\\
无粘激波–气泡，$|\phi_0|<0.08$ & 0.07 & 1.0 & 3.6\\
无粘激波–气泡，$|\phi_0|<0.04$ & 0.11 & 1.1 & 4.6\\
粘性激波–气泡（$\Delta x=0.005$），$|\phi_0|<$0.04 & 0.044 & 0.8 & 0.5\\
粘性激波–气泡（$\Delta x=0.005$），$|\phi_0|<$0.08 & 0.027 & 0.5 & 0.4\\
粘性激波–气泡（$\Delta x=0.005$），$|\phi_0|<$0.15 & 0.033 & 0.5 & 0.4\\
激波–等熵涡，涡核 $r_0<2R_v$ & 0.039 & 0.2 & 0.7\\
激波–等熵涡，涡核 $r_0<R_v$ & 0.254 & 0.4 & 0.5\\
单模 RM，$|\phi_0|<$0.08 & 0.029 & 2.2 & 0.2\\
单模 RM，$|\phi_0|<$0.03 & 0.032 & 5.4 & 1.8\\
\bottomrule\end{tabularx}\end{table}
在四个 DNS 中，位置形式与线积分的中位相差都在 0.25\% 以内，而旧的面积求积在卷起后偏离可达数十个百分点。这说明位置形式给出的是精确的边界环量，此前的差异来自参照本身。需要说明，示踪点是用同一速度场积分的，所以这项检验确认的是“位置形式环量 = 同一回路上的速度环量”；环量的物理正确性仍依赖 DNS。

\section{扩散：方程与密度估计的分解}
\begin{table}[!htbp]\centering\small\caption{末时刻扩散误差（RMS 误差除以参照标准差）。}
\begin{tabularx}{0.8\linewidth}{@{}L r r@{}}\toprule
变体 & Sod & 爆炸波\\\midrule
精确 $\rho$（oracle） & 0.163\% & 1.85\%\\
账本 $\rho$（论文默认） & 0.174\% & 2.19\%\\
SPH $\rho$ & 0.195\% & 2.48\%\\
DTFE $\rho$ & 0.253\% & 9.12\%\\
面积度量 & 1.846\% & 22.30\%\\
账本 $\rho$，每帧 4 个子步 & 0.151\% & 2.13\%\\
账本 $\rho$，Crank–Nicolson & 0.144\% & 2.29\%\\
账本 $\rho$，边电导 M 矩阵 & 1.773\% & 2.15\%\\
\bottomrule\end{tabularx}\end{table}
读数如下。
\begin{itemize}
\item 质量度量与面积度量相差 10 倍以上，而电导里换用不同的密度估计只差 1.1–1.5 倍（DTFE 在爆炸波中为 4.9 倍）。所以“4–10 倍”检验的是方程的选择，账本贡献的是让这一选择只凭位置就能实现。
\item 子步与 Crank–Nicolson 对结果影响很小，快照间隔不是误差的主要来源。
\item 改用气体动理学电导 $\rho D\propto T^{0.7}$，$T$ 由位置重建（沿示踪点熵守恒，检出激波时按实测密度比给出 Rankine–Hugoniot 熵跃；$T$ 的中位误差为机器精度，95 分位 2.8\%）。结果：质量度量 0.126\%（用精确 $T$ 为 0.126\%），面积度量 2.9\%，差距依然存在。
\end{itemize}

\section{账本 v3：多窗口检测、自动推迟、张量/混合估计}
\paragraph{改动。}
\begin{itemize}
\item \textbf{多窗口激波检测}：一个星要经过多帧才被激波扫过时（$C_s\lesssim1$、弱激波或粘性激波），在长度 $w\le6$ 帧的窗口内，若星面积下降超过 $\max(\tau,3w\,\mathrm{med}_t|R_i|)$，并且该窗口上的最小二乘仿射映射有一个主伸长等于 1（容差 $0.1$，即签名的转动不变形式），就把窗口内各步标为热步。第一版没有签名确认，卷起区被误标，激波–气泡的 95 分位从 29.7\% 升到 39.4\%；加确认后为 30.3\%。
\item \textbf{自动推迟上限}：取热步持续长度 99 分位的 2 倍，限制在 20–50 帧之间。
\item \textbf{张量账本}：对链式最小二乘变形梯度（12 近邻二次拟合，推迟规则相同）取 $\ln\det$。光滑流中是二阶，跨平面激波精确，但偶有离群值。
\item \textbf{混合估计}：张量值与多边形值在 $\ln J$ 上相差不超过 $0.1$ 时用张量值（86--100\% 的示踪点），否则用多边形值；质量仍由多边形账本给出。
\end{itemize}
\begin{table}[!htbp]\centering\small\caption{末时刻密度误差（中位 / 95 分位）：旧账本与混合估计，以及指示量识别最差 5\% 误差的 AUC。激波–等熵涡取 $t=0.6$，粘性激波–气泡取 $\Delta x=0.005$。}
\begin{tabularx}{\linewidth}{@{}L r r r@{}}\toprule
算例 & 旧账本（v2） & 混合估计（v3） & 指示量 AUC\\\midrule
光滑制造映射 & 0.04\% / 0.1\% & 0.01\% / 0.0\% & 0.53\\
会聚柱面制造激波 & 0.15\% / 0.7\% & 0.05\% / 0.4\% & 0.89\\
弱激波（$\eta=1.5,C_s=0.5$） & 0.36\% / 5.8\% & 0.00\% / 0.0\% & 0.59\\
慢激波（$C_s=0.25$） & 7.58\% / 25.2\% & 0.00\% / 0.0\% & 0.50\\
Sod & 0.00\% / 1.0\% & 0.00\% / 0.3\% & 0.97\\
柱面爆炸波 & 0.58\% / 32.2\% & 0.29\% / 33.4\% & 0.91\\
双环流 & 1.24\% / 5.6\% & 0.11\% / 0.5\% & 0.74\\
激波–气泡（无粘） & 0.67\% / 29.7\% & 0.40\% / 31.0\% & 0.93\\
激波–等熵涡 & 0.34\% / 3.4\% & 0.20\% / 1.9\% & 0.94\\
单模 RM & 0.73\% / 35.4\% & 0.47\% / 35.3\% & 0.90\\
激波–气泡（粘性） & 0.53\% / 8.9\% & 0.31\% / 9.3\% & 0.86\\
\bottomrule\end{tabularx}\end{table}
\begin{table}[!htbp]\centering\small\caption{0.02$h$ 位置噪声下的末时刻密度误差（中位 / 95 分位）。}
\begin{tabularx}{0.85\linewidth}{@{}L r r@{}}\toprule
算例 & 多边形账本 & 混合估计\\\midrule
光滑制造映射 & 1.30\% / 3.9\% & 0.82\% / 2.5\%\\
Sod & 1.45\% / 5.2\% & 0.94\% / 3.3\%\\
柱面爆炸波 & 3.30\% / 32.8\% & 2.25\% / 34.0\%\\
激波–气泡（无粘） & 3.12\% / 30.9\% & 2.27\% / 30.6\%\\
激波–等熵涡 & 1.94\% / 7.8\% & 1.22\% / 5.7\%\\
单模 RM & 3.70\% / 35.7\% & 2.66\% / 35.6\%\\
激波–气泡（粘性） & 2.85\% / 12.0\% & 2.05\% / 12.0\%\\
\bottomrule\end{tabularx}\end{table}
\paragraph{读数。}
\begin{itemize}
\item 重跑的三个 DNS 都复现了旧数值：激波–涡 0.34\%/3.4\%，RM 0.73\%/35\%，粘性激波–气泡 0.53\%/8.9\%。
\item 除 Sod（中位误差本已是机器精度，95 分位从 0.97\% 降到 0.29\%）外，混合估计在所有流动算例中都降低了中位误差。激波–等熵涡的 95 分位也明显下降（涡核内的非仿射误差正是张量账本能纠正的）；激波、接触间断与界面附近的 95 分位基本不变。
\item 噪声下混合估计的中位误差降低 27\%–37\%。
\item 指示量在激波流中能很好地识别最大误差（AUC 0.86–0.97），在光滑流中没有信息，它不是误差上界。
\end{itemize}

\section{仍未解决}
\begin{enumerate}
\item 张量账本有离群值，目前靠混合规则兜底；激波处的 95 分位没有改善。
\item 多窗口检测只在制造激波上验证过，尚未在可分辨的粘性激波（SBLI 类）上检验。
\item T4 的冲击/光滑拆分还没有改在合并块上定义，也没有与 DNS 斜压力矩直接比较。
\item 粘性激波–气泡这次只重跑了 $\Delta x=0.005$；三维与真实 PTV 数据仍未做。
\end{enumerate}

\clearpage\part{附录部分　后续工作：轨迹网络上的脊分类（双曲 LCS、剪切型脊、激波）}


\noindent\fcolorbox{accent!40}{tint}{\begin{minipage}{0.965\linewidth}\vspace{0.3em}
\textbf{\sffamily 问题}\quad AFRC 论文发现 \AFH{} 与 FTLE 的脊相符。但 FTLE 的局部极大值不一定是双曲 LCS：它也可能是没有法向排斥的剪切型脊（Haller 2011 的平行剪切反例），在可压缩流中还可能是激波。本稿给出只凭示踪点位置区分这三者（以及光滑波）的判据与初步检验。它是 T1–T9 之后的工作，用到 T1 的热步与 T2 的账本，但不属于 T1–T9。

\textbf{\sffamily 结论}\quad
(1) 只看窗口两端的变形梯度，原则上分不开双曲与剪切。(2) 伸长边的增长指数 $\beta$（距离对时间的双对数斜率）可以分开二者：剪切 $\beta\to1$，定常双曲 $\beta=\lambda T/(2\ln2)$；鞍点叠加剪切时给出了半窗比的渐近式。(3) 激波用跳跃占比 $K$ 识别，光滑波用面积占比 $W$ 识别。(4) 在剪切层上判别 100\% 正确；后向 FTLE 落在激波上的极大值中，判为激波的点 99.8\%–100\% 正确。(5) 在非定常、间歇拉伸的双环流中，单窗口的 $\beta$ 与精确判别只有 60\% 一致；有限点对使 $\beta$ 偏高，加密后偏差下降但较慢。所以目前只在定常或统计平稳的窗口内可靠。
\vspace{0.3em}\end{minipage}}




\section{两端变形梯度不可辨识}

\begin{proposition}\status{已证}
取 $\gamma T=\sigma-\sigma^{-1}$，鞍型 $F_h=\mathrm{diag}(\sigma,\sigma^{-1})$ 与简单剪切 $F_s=I+\gamma T\,\mathbf e_1\mathbf e_2^\top$ 的奇异值相同，二者只差一个左转动（极分解中剪切的转角满足 $\tan\varphi_R=\gamma T/2$）。因此任何只依赖 $C=F^\top F$ 的量，包括 FTLE、点对拉伸大小与 \AFH{} 的边权，在单点上都不能区分二者。
\end{proposition}

激波的不同之处在于 $J\neq1$，而且变化发生在一步之内；稀疏扇或波后持续膨胀则是光滑的面积变化。所以需要的信息是：拉伸随时间如何增长，以及拉伸伴随的面积变化是跳跃的还是光滑的。这些都在网络的边时间序列与 T2 账本中。

\section{指标}

对每个候选节点（脊场取 90\% 分位以上），取其\textbf{伸长边} $e^\ast(i)=\arg\max_{e\ni i}[\ell_e(T)-\ell_e(0)]$，$\ell_e=\ln|q_e|$。热步由 T1 的 C 通道给出，$\ell^c$ 只累加非热步的增量。

\paragraph{增长指数。}
\begin{equation}
\beta_i=\frac{\ell^c_{e^\ast}(T)-\ell^c_{e^\ast}(T/2)}{\ln2},\qquad R_i=\frac{\ell^c_{e^\ast}(T)-\ell^c_{e^\ast}(T/2)}{\ell^c_{e^\ast}(T/2)-\ell^c_{e^\ast}(0)}.
\end{equation}

\begin{proposition}[解析值]\status{已证+数值}
(i) 简单剪切：$|F\hat r|=\sqrt{(\cos\varphi+\gamma t\sin\varphi)^2+\sin^2\varphi}$，$\gamma T\gg1$ 时 $\beta\to1$，$R\to\ln2/\ln(\gamma T/2)$。
(ii) 定常鞍点：$\beta=\lambda T/(2\ln2)$，$R=1$。
(iii) 鞍点叠加剪切 $u=(\lambda x+\gamma y,-\lambda y)$：$F(t)=\begin{psmallmatrix}e^{\lambda t}&(\gamma/\lambda)\sinh\lambda t\\0&e^{-\lambda t}\end{psmallmatrix}$，当 $e^{-\lambda T/2}\ll1$ 时
\begin{equation}
R\approx\frac{\lambda T}{\lambda T+\ln(1+\gamma^2/4\lambda^2)} .
\end{equation}
\end{proposition}
\begin{proof}[推导要点]
(iii) $t\gg\lambda^{-1}$ 时 $|F\hat r|\approx e^{\lambda t}|\cos\varphi+\tfrac{\gamma}{2\lambda}\sin\varphi|$；取伸长最大的方向，$\ln|F\hat r|\approx\lambda t+\tfrac12\ln(1+\gamma^2/4\lambda^2)$，代入 $R$ 的定义即得。
\end{proof}

式 (iii) 说明，一个弱鞍点叠加强剪切，在有限窗口内与纯剪切无法区分：$R>\tfrac12$ 等价于 $\lambda T>\ln(1+\gamma^2/4\lambda^2)$，即指数拉伸要超过剪切贡献的对数拉伸。这是窗口有限造成的信息极限。

\paragraph{跳跃占比与面积占比。}
$K_i=\sum_{\text{热步}}|\Delta\ell_{e^\ast}|/\sum|\Delta\ell_{e^\ast}|$，激波穿越使 $K\to1$。
$W_i=|\Delta\ln J^c_i|/|\Delta\ell^c_{e^\ast}|$，$\ln J$ 取自 T2 账本：沿边方向的单轴膨胀或压缩使 $W\approx1$，不可压的双曲或剪切中 $W\approx0$。

\paragraph{判别规则。}$K>0.5$ 为激波；否则 $W>0.5$ 为光滑波；否则 $\beta>1.5$ 为双曲，$0.5\le\beta\le1.5$ 为剪切，$\beta<0.5$ 为“已停止”（拉伸发生在窗口前半）。剪切的判区以 $\beta=1$ 为中心；定常双曲要 $\beta>1.5$ 需 $\lambda T>2.1$，这也是形成可见脊所需的拉伸量。

\paragraph{两个不采用的做法。}
\begin{itemize}
\item \textbf{星内“拉伸接力”秩相关 $\zeta$。}纯鞍点为 0.997，但叠加 $\gamma/\lambda=0.5$ 的弱剪切就降到 0.13，不稳健。
\item \textbf{物质坐标系的 QR 分解} $\ln|Fn_0|=\ln J-\ln\lambda_t+\tfrac12\ln(1+\cot^2\chi)$（数值验证到 $7\times10^{-16}$）。它给出 Haller 意义下法向排斥的严格定义，但要求脊的初始切向误差小于 $\sigma_1^{-1}$，在充分发展的拉伸区无法满足。只作为定义使用。
\end{itemize}

\begin{figure}[!htbp]\centering
\includegraphics[width=\linewidth]{fig1_theory.pdf}
\caption{(a) 鞍点叠加剪切的半窗比 $R$：实线为 24 方向星上的数值，虚线为渐近式（$\lambda T=3$、$\gamma/\lambda=5$ 时为 0.63 对 0.60）。(b) 增长指数：纯剪切趋于 1（蓝带为剪切判区），同 $\sigma_1$ 的鞍点为 $\lambda T/(2\ln2)$。(c) 激波块的签名，黑线为单轴曲线（T1 的内容，此处用于说明激波类）。}\label{pE:fig:th}
\end{figure}

\section{检验}

\paragraph{设置。}候选点取 \AFH{}（单尺度、原门控、$k=24$）或位置法 FTLE 的 90\% 分位以上。精确 $\beta$ 由 $\pm10^{-6}$ 扰动点的 RK4 积分得到 $F(t)$，取 $\beta^{\rm ex}=[\ln\sigma_1(T)-\ln\sigma_1(T/2)]/\ln2$。

\subsection{基准流}
\begin{table}[!htbp]\centering\small
\caption{基准流上 \AFH{} 候选点的分类。“一致率”为与精确 $\beta$ 判别的一致比例。}\label{pE:tab:bench}
\begin{tabularx}{\linewidth}{@{}l r r r r r r L@{}}\toprule
算例 & 候选数 & 双曲 & 剪切 & 已停止 & 一致率 & Spearman & 网络 $\beta$ 中位数（精确）\\\midrule
平行剪切层 $u=\tanh(y/0.1)$ & 413 & 0\% & 100\% & 0\% & 1.00 & 0.54$^\dagger$ & 0.996（0.993）\\
局部鞍点 & 503 & 76\% & 13\% & 11\% & 0.60 & 0.975 & 2.46（1.47）\\
鞍点+剪切层 $\gamma=2$ & 641 & 86\% & 10\% & 3\% & 0.55 & 0.89 & 2.30（1.55）\\
鞍点+剪切层 $\gamma=10$ & 710 & 11\% & 47\% & 42\% & 0.80 & 0.58 & 0.51（0.49）\\\bottomrule
\end{tabularx}\par\smallskip{\footnotesize $^\dagger$ $\beta$ 全部集中在 0.95–1.00，秩相关不说明问题。}
\end{table}

剪切层（Haller 的反例：FTLE 有脊、没有法向排斥）全部判为剪切。局部鞍点上网络 $\beta$ 与精确值秩相关 0.975，但系统偏高：网络读的是跨过稳定流形的有限点对，两端分到两侧后持续分离；精确 $\beta$ 读的是无穷小邻域，离开鞍点区后拉伸就停止。就“这条脊是不是双曲结构”而言，有限点对给出的 76\% 更接近构造上的真值，但这也意味着“精确 $\beta$”本身不是理想的参照。

\begin{figure}[!htbp]\centering
\includegraphics[width=\linewidth]{fig2_bench.pdf}
\caption{基准流。左：\AFH{} 场（初始位置）；中：候选点分类；右：网络 $\beta$ 与精确 $\beta$。}
\end{figure}

\subsection{双环流与 Lamb–Oseen 涡对}
\begin{table}[!htbp]\centering\small
\caption{不可压算例。}
\begin{tabularx}{\linewidth}{@{}L r r@{}}\toprule
& 双环流（$T=10$） & Lamb–Oseen（$N=8000$）\\\midrule
网络：双曲 / 剪切 / 已停止 & 58\% / 35\% / 7\% & 78\% / 21\% / 1\%\\
精确 $\beta$：双曲 / 剪切 / 已停止 & 41\% / 39\% / 20\% & 50\% / 36\% / 14\%\\
一致率；$\beta$ 的 Spearman & 0.60；0.51 & 0.63；0.74\\\bottomrule
\end{tabularx}\end{table}

Lamb–Oseen 的 \AFH{} 候选点落在共转参考系的 8 字形分界线上，即共转系鞍点的稳定、不稳定流形，网络判为以双曲为主。双环流中约三分之一的候选点判为剪切。原因在理论上是预期的：一个强迫周期内，示踪点只在经过鞍点区时被拉伸，拉伸是间歇的，$\beta$ 所假设的“窗口内平稳”不成立。所以在非定常流动中，$\beta$ 判为剪切只表示“窗口后半拉伸放缓”，不能直接解释为剪切 LCS。

Lamb–Oseen 的分辨率阶梯表明有限点对偏差在收敛：
\begin{center}\small
\begin{tabular}{lccc}\toprule
$N$ & 2000 & 8000 & 32000\\\midrule
网络 $\beta$ − 精确 $\beta$（中位数） & 0.50 & 0.45 & 0.26\\
与精确判别的一致率 & 0.59 & 0.63 & 0.74\\\bottomrule
\end{tabular}\end{center}

\begin{figure}[!htbp]\centering
\includegraphics[width=\linewidth]{fig3_dg_lo.pdf}
\caption{上：双环流；下：Lamb–Oseen 涡对。(1) 只凭位置的步涡量（T3）；(2) \AFH{}；(3) 候选点分类；(4) 网络 $\beta$ 与精确 $\beta$。}
\end{figure}

\subsection{可压缩算例}
\begin{table}[!htbp]\centering\small
\caption{可压缩算例的候选点分类。}
\begin{tabularx}{\linewidth}{@{}L r r r@{}}\toprule
& Sod & 爆炸波 & 激波–等熵涡（DNS）\\\midrule
\AFH{} 候选点的位置 & 全在稀疏扇 & 被激波扫过、随后膨胀的壳层 & 涡周围\\
\AFH{} 候选点分类 & 光滑波 100\% & 光滑波 97\% & 双曲 46\% / 剪切 42\% / 已停止 12\%\\
后向 FTLE 候选点中被激波扫过的比例 & 100\% & 99.6\% & —（全部示踪点都被扫过）\\
后向 FTLE 判为激波的比例；其中正确的比例 & 84\%；100\% & 64\%；99.8\% & 0\%\\
后向 FTLE 其余分类 & 剪切 7\%、已停止 9\% & 已停止 35\% & 剪切 99\%\\
前向 FTLE 候选点分类 & 光滑波 100\% & 光滑波 100\% & 双曲 50\% / 剪切 50\%\\\bottomrule
\end{tabularx}\end{table}

三点观察。
\begin{enumerate}
\item \textbf{\AFH{} 本身几乎不在激波上产生候选点}，即使关掉门控也是如此：它的边权是增量的中位数，单次跳跃不改变中位数。Sod 与爆炸波中的 \AFH{} 脊来自稀疏扇与波后持续膨胀，$W$ 把它们判为光滑波，而不是双曲结构。所以“分离激波”的需求实际上来自 FTLE 一类的诊断。
\item \textbf{FTLE 落在激波上的极大值可以识别。}后向 FTLE 的候选点中，判为激波的点 99.8\%–100\% 正确；没有判为激波的（16\%–36\%）跳跃占比在 0.3–0.5，原因是 C 通道漏掉了部分穿越步，或波后膨胀摊薄了跳跃。
\item \textbf{激波–等熵涡中 FTLE 的极大值来自涡，而不是激波。}$M_s=1.3$ 的激波较弱（$\ln\eta\approx0.4$），后向 FTLE 的极大值 99\% 判为剪切，来自涡的差动旋转，与物理图像一致。
\end{enumerate}

\begin{figure}[!htbp]\centering
\includegraphics[width=\linewidth]{figL4_comp.pdf}
\caption{后向 FTLE 候选点的分类：(a) Sod，(b) 爆炸波，(c) 激波–等熵涡。}
\end{figure}

\section{现状与下一步}

\paragraph{已经成立的部分。}两端 $F$ 不可辨识的命题；$\beta$、$R$ 的解析值与鞍点叠加剪切的渐近式；激波（$K$）与光滑波（$W$）的识别；剪切层上的完全正确判别。

\paragraph{尚未解决的部分。}
\begin{enumerate}
\item \textbf{非平稳拉伸。}单窗口的 $\beta$ 在间歇拉伸的流动中不可靠（双环流一致率 0.60）。建议改为多窗口统计：在一系列嵌套子窗口 $[0,T/2^k]$ 上拟合 $\ell(t)$ 的增长律，或在滑动子窗口内逐段判别再按时间加权。
\item \textbf{参照本身。}点态的精确 $\beta$ 在局部鞍点中把离开鞍点区的点判为“已停止”，而有限点对把整条脊判为双曲。需要一个面向“脊”的参照，例如沿脊曲线的 Haller 法向排斥率，在可解析的流动上计算。
\item \textbf{有限点对偏差}随加密下降但慢于一阶（0.50→0.26，$N$ 增大 16 倍）。
\item \textbf{判别阈值}（$K,W$ 取 0.5，$\beta$ 取 0.5 与 1.5）来自解析值，没有做敏感性分析。
\end{enumerate}

\paragraph{建议的检验顺序。}(1) 多窗口 $\beta$ 在双环流上的一致率；(2) 平面斜激波与正激波，检验 $K$ 与倾角无关；(3) Re = 1000 的激波—边界层相互作用：滑移线应判为剪切（AFRC 论文中它表现为一对紧邻的排斥、吸引脊），$\lambda$ 激波脚附近 FTLE 的竖直带应判为激波，卷起之间的 braid 应判为双曲。



\clearpage\part{总结}
\section{最终结论}
\begin{enumerate}
\item \textbf{理论统一。}成熟部分是质量度量下同一个 Dirac 算子的不同块：复上链的实部与虚部给出环量与账本，$D^2$ 的节点块给出变密度扩散，其时间平均给出加权动态 Laplacian；三点映射的合成给出激波签名与链式张量。所有量都是位置的有限步代数恒等式，不需要速度、不依赖步长，跨激波时用合并块或推迟重锚仍然成立。
\item \textbf{度量即测度。}在可压缩流中，以质量而不是面积为度量是可检验的物理选择：扩散误差相差 3.8--10.6 倍（来自方程，在 $\rho D\propto T^{0.7}$ 下依然存在），组分按计数偏差 2.5 倍。
\item \textbf{精度与稳健性。}只凭位置的密度比 DTFE/SPH 准 1–2 个数量级；噪声误差 $\approx0.74\,\sigma/h$ 且不随累积变形放大；稀疏与断轨下可用（逐轨账本覆盖率 100\%）。
\item \textbf{物理结果。}激波–重气泡界面在激波扫过后先被切向缩短、抵消大部分压缩；单模 RM 到 $ka\approx0.6$ 时界面梯度增长约 94\% 来自压缩，增长率与 Richtmyer 模型相差 3\%；只凭位置的沉积环量与 Samtaney–Zabusky 模型相差 8\%。
\item \textbf{验证。}全部代数恒等式通过精确符号验证；账本在制造流映射上平面激波处精确、弯曲激波处误差随 $h/R_s$、光滑流中为一阶（张量账本为二阶）；环量在同一物质边界上与 DNS 线积分一致到 0.25\% 以内。\item \textbf{边界。}卷起后界面带的结论依赖雷诺数与示踪点分辨率（信息极限）；爆炸波薄壳中断轨后的新轨迹标定误差大；三维与真实 PTV 数据尚未做。
\end{enumerate}

\section{论文计划}
\begin{itemize}
\item \textbf{论文 A}（本报告的主体）：“质量度量下的拉格朗日 Hodge–Dirac 计算：只凭示踪点位置的可压缩输运与掺混”，包含 T1、T2、T3（局部）、T4、T6、T7（动态 Laplacian）、T8、T9 与稳健性；英文初稿见另附的论文 PDF 与 TeX。
\item \textbf{拓扑信号文章}：T3 的全局分解、T5 的谱与声学 Dirac、DESP，以及以质量度量 Laplacian 定义的 von Neumann 熵与多路网络（Bianconi 方向）。
\item \textbf{脊分类文章}：附录部分的双曲 / 剪切 / 激波判别，需先解决非平稳拉伸与有限点对偏差。
\end{itemize}

\section{代码与 notebook}
代码包 \texttt{LHD\_code.zip} 的结构：\texttt{lhd/}（核心模块：\texttt{lhdp.py} 三点映射与有限步恒等式，\texttt{lhdx.py}/\texttt{lhd.py} 复形与账本，\texttt{bcom.py} 账本变体、SPH 传导与误差统计，\texttt{stream.py} 逐轨账本，\texttt{netdiff.py} 网络扩散，\texttt{stretch\_est.py} 链式张量与 Nanson 余向量，\texttt{classify.py} 与 \texttt{afrc\_h\_engine.py} 脊分类，\texttt{cases.py} 解析与精确算例，\texttt{euler2d.py}/\texttt{euler2d\_visc.py}/\texttt{sv\_case.py} DNS）；\texttt{scripts/}（每节结果的脚本）；\texttt{results/}（全部 \texttt{res\_*.json}）；\texttt{LHD\_demo.ipynb}（可运行演示：有限步恒等式、Sod 账本与激波签名、质量形式扩散、爆炸波掺混归因、动态 Laplacian、断轨账本）。v2 新增：\texttt{lhd/mfm.py}（制造流映射）、\texttt{lhd/ledger\_ub.py}（指示量、多窗口检测、\texttt{ledger\_v3}、张量/混合账本、热力学重建）、\texttt{tests/}（精确符号检验与变形关系检验）、\texttt{tools/}（数字单一来源与核对）、\texttt{scripts/v1\_*}–\texttt{v7\_*}（验证脚本）。

\end{document}